\documentclass[../main.tex]{subfiles}
\begin{document}
\ifSubfilesClassLoaded{\section{Optimal Transport (OT)}}

\subsection{Transport Maps}
\begin{definition}
    A Polish space is a completely metrizable and separable topological space $(X,\tau)$. 
    The Borel $\sigma$-algebra of $X$ is the $\sigma$-algebra $\mathcal{B}(X):=\sigma(\tau)$ generated by the topology of the space.
\end{definition}

In the following, we will always consider probability measures $\mu$ in some Polish space $X$, and we will denote the space of all probability measures in $X$ with the symbol $\mathcal{P}(X)$.

\begin{definition}
    The pushforward by a measurable map $T: X \to Y$ is the induced map $T_\#: \mathcal{P}(X) \to \mathcal{P}(Y)$ such that $T_\#\mu = \mu \circ T^{-1}$, or equivalently $(T_\#\mu)(B) := \mu(T^{-1}(B))$ for all $B\in\mathcal{B}(Y)$.
\end{definition}

\begin{definition}
    A transport map between $\mu\in\mathcal{P}(X)$ and $\nu\in\mathcal{P}(Y)$ is a measurable map $T: X \to Y$ such that $T\mu = \nu$.
\end{definition}

We will denote the space of the transport maps by $\mathcal{T}(\mu, \nu) := \{T: X \to Y \mid T\mu = \nu\}$. \\
Intuitively, a transport map between two distributions $\mu$ and $\nu$ moves points' mass from the first distribution to the second one.

\begin{example}
\label{ex:dirac_tranport}
    Let $\delta_1, \delta_2\in\mathcal{P}(\R)$ be Dirac measures on $\R$, and $T: \R \to \R$ such that $T(x)=2x$. Then, we have $T\delta_1(B) = \delta_1(T^{-1}(B)) = \delta_2(B)$ since $T(1) = 2$. Hence, $T\in\mathcal{T}(\delta_1, \delta_2)$ is a transport map between $\delta_1$ and  $\delta_2$. Intuitively, the map $T$ transports all the mass of $\delta_1$, which is concentrated at the point $1$, to the point $2$. \\
    More generally, Dirac measures can only be transported to other Dirac measures by any measurable map $T: X \to Y$. Indeed, consider the measure $\delta_x\in\mathcal{P}(X)$, and a measurable map $T: X \to Y$. Then, for all $B\in\mathcal{B}(Y)$, we have 
    $$T\delta_x(B) = \delta_x(T^{-1}(B)) = \delta_{T(x)}(B).$$
    Note that the other mappings $x \neq y \mapsto T(y)$ are not relevant, as they have zero mass.
\end{example}

\begin{remark}
    In general, $\mathcal{T}(\mu, \nu) = \varnothing$ given $\mu\in\mathcal{P}(X)$, $\nu\in\mathcal{P}(Y)$.
    For instance, there exists no transport maps from a Dirac measure $\delta_x\in\mathcal{P}(X)$ to a non Dirac measure $\nu$.
\end{remark}
\vspace{10pt}

\begin{example} \hfill
\label{ex:real_transports}
    \begin{enumerate}[(i)]
        \item Let $\nu\in\mathcal{P}(\R)$ be a probability measure on the real line $\R$. The quantile function of $\nu$ is the map 
        $$
        q_\nu : [0,1]\to \R \text{ such that } q_\nu(u) := \inf \{x\in\R \mid F_\nu(x) \geq u\}.
        $$
        Then, $q_\nu\in\mathcal{T}(\lambda, \nu)$ is a transport map from the Lebesgue measure $\lambda$ on $[0,1]$ to the measure $\nu$. Indeed, for any random variable $X$, the variable $q_X(U)\sim X$ has the same law of $X$, where $U\sim U(0,1)$ is a uniform variable on $[0,1]$.
        \item Let $\mu\in\mathcal{P}(\R)$ be a continuous probability measure on the real line $\R$. The cumulative function of $\mu$ is the map
        $$
        F_\mu : \R \to [0,1] \text{ such that } F_\mu(x) := \mu((-\infty, x]).
        $$
        Then, $F_\mu\in\mathcal{T}(\mu, \lambda)$ is a transport map from $\mu$ to the Lebesgue measure $\lambda$. Indeed, for any continuous random variable $X$, the variable $F_X(X)\sim U$ has the same law of a uniform distribution $U\sim U(0,1)$. 
    \end{enumerate}
\end{example}

\begin{lemma}
\label{lemma:transports_composition}
    Let $T_1: X \to Y$ be a transport map from $\mu\in\mathcal{P}(X)$ to $\nu\in\mathcal{P}(Y)$, and $T_2: Y \to Z$ be a transport map from $\nu\in\mathcal{P}(Y)$ to $\rho\in\mathcal{P}(Z)$. Then, $T_2\circ T_1 : X \to Z$ is a transport map
    $$
    T_2 \circ T_1 \in\mathcal{T}(\mu, \rho)
    $$
\end{lemma}
\begin{proof}
    Let $B\in\mathcal{B}(Z)$, then we have
    \begin{align*}
    ((T_2\circ T_1)_\#\mu)(B) 
    &= \mu((T_2\circ T_1)^{-1}(B)) 
    = \mu(T^{-1}_1(T^{-1}_2(B))) = \\
    &= \nu(T^{-1}_2(B)) =  \text{ since } T_1\in\mathcal{T}(\mu, \nu)\\
    &= \rho(B) \text{ since } T_2\in\mathcal{T}(\nu, \rho).
    \end{align*}
\end{proof}

\begin{proposition}
    Let $\mu, \nu\in\mathcal{P}(\R)$ be probability measures with $\mu$ continuous. Then, the map $q_\nu \circ F_\mu : \R \to \R$, called Monge map, is a transport map from $\mu$ to $\nu$, i.e.
    $$
    q_\nu \circ F_\mu \in \mathcal{T}(\mu, \nu).
    $$
\end{proposition}
\begin{proof}
    Recall that $F_\mu\in\mathcal{T}(\mu, \lambda)$ and $q_\nu\in\mathcal{T}(\lambda, \nu)$ by (\ref{ex:real_transports}), where $\lambda$ is the Lebesgue measure on $[0,1]$. Then, the thesis follows from (\ref{lemma:transports_composition}).
\end{proof}
\noindent
Intuitively, the Monge map $q_\nu\circ F_\mu$ transports mass from $\mu$ to $\nu$ in a monotone way.\\

In the setting $X = Y = \R^d$, it is particularly easy to prove if a measurable map $T: \R^d \to \R^d$ is a transport map between two given distributions $\mu, \nu\in\mathcal{P}(\R^d)$, as shown by the following proposition.
\begin{proposition} (Monge - Amp\'ere) \ 
    $\mu, \nu\in\mathcal{P}(\R^d)$ absolutely continuous respect to $\lambda$ Lebesgue measure, $T: \R^d \to \R^d$ injective, $C^1$ and $\mid JT \mid \neq 0$. Then
    \begin{equation}
    \label{eq:Monge-Ampere}
    T\in\mathcal{T}(\mu, \nu) \iff T \text{ sol.\ of the Monge-Amp\'ere eq.\ } \det(JT)\cdot \left(\frac{d\nu}{d\lambda} \circ T\right) = \frac{d\mu}{d\lambda}.
    \end{equation}
\end{proposition}
\begin{proof}
    Suppose that $T$ is a solution of $(\ref{eq:Monge-Ampere})$, and we prove that $\mu(T^{-1}(B)=\nu(B)$ for all $B\in\mathcal{B}(\R^d)$.
    Then, it holds, 
    \begin{align*}
        \mu(T^{-1}(B)) &= \int_{T^{-1}(B)} \frac{d\mu}{d\lambda}(\underline{x}) d\underline{x} = \text{ since } \mu \ll\lambda \\
        &= \int_{T^{-1}(B)} \det(JT(\underline{x}))\frac{d\nu}{d\lambda}(T(\underline{x})) d\underline{x} = \\
        &= \int_B \frac{d\nu}{d\nu}(\underline{y}) d\underline{y}= \quad \text{by change of variable } \underline{y} = T(\underline{x}) \\ 
        &= \nu(B).
    \end{align*}
    Conversely, suppose that $T$ is a transport from $\mu$ to $\nu$, and we show that $T$ solves the Monge - Amp\'ere equation. Thus, we have for all $B\in\mathcal{B}(\R^d)$
    \begin{align*}
        \nu(B) &= \int_B \frac{d\nu}{d\lambda}(\underline{y})d\underline{y} = \int_{T^{-1}(B)} \det(JT(\underline{x}))\frac{d\nu}{d\lambda}(T(\underline{x})) d\underline{x} \\
        \mu(T^{-1}(B)) &= \int_{T^{-1}(B)} \frac{d\mu}{d\lambda}(\underline{x}) d\underline{x}.
    \end{align*}
    Hence, we have for all $B\in\mathcal{B}(\R^d)$
    $$
    \int_{T^{-1}(B)} \det(JT(\underline{x}))\frac{d\nu}{d\lambda}(T(\underline{x})) d\underline{x}
    = \int_{T^{-1}(B)} \frac{d\mu}{d\lambda}(\underline{x}) d\underline{x},
    $$
    which implies that $T$ solves (\ref{eq:Monge-Ampere}) by fundamental theorem of variational calculus.
\end{proof}

\subsection{Couplings}
We introduce a generalized notion of transport maps which will be useful to define optimal transport problems.

\begin{definition}
    A coupling of two probabilities $\mu\in\mathcal{P}(X), \nu\in\mathcal{P}(Y)$ is a probability measure $\pi\in\mathcal{P}(X\times Y)$ with marginals $\mu, \nu$, i.e. such that $\pi(A\times Y) = \mu(A)$ and $\pi(X\times B)=\nu(B)$ for all $A\in\mathcal{B}(X), B\in\mathcal{B}(Y)$. 
\end{definition}

We will denote the space of couplings between $\mu$ and $\nu$ by $Cpl(\mu,  \nu)$.\\
Equivalently, we can define a coupling between $\mu$ and $\nu$ as a probability measure $\pi\in\mathcal{P}(X\times Y)$ which induces $\mu, \nu$ by pushforward under projections $p_X : X\times Y \to X$ and $p_Y : X\times Y \to Y$, i.e. such that $p_X\pi = \mu$ and $p_Y\pi = \nu$. Indeed, for $A\in\mathcal{B}(X), B\in\mathcal{B}(Y)$ we have
\begin{align*}
    (p_X\pi)(A) &= \pi(p^{-1}_X(A)) = \pi(A\times Y) = \mu(A) \\
    (p_Y\pi)(B) &= \pi(p^{-1}_Y(B)) = \pi(X\times B) = \nu(B).
\end{align*}

A further equivalent characterization of the coupling property is given by the following proposition.

\begin{proposition}
\label{prop:characterization_cpl}
    Let $\mu, \nu$ be probability measures on $X, Y$ respectively. Then,
    $$
    Cpl(\mu, \nu) = 
    \left\{\pi\in\mathcal{P}(X\times Y) \mid \int f d\pi = \int f\mu,
    \int g d\pi = \int g\nu, \ f\in L^1(\mu), \ g\in L^1(\nu)\right\}
    $$
\end{proposition}
\begin{proof}
    Let $\pi\in\mathcal{P}(X\times Y)$ such that $\int_{X\times Y} f d\pi = \int_X f\mu$ and $ \int_{X\times Y} g d\pi = \int_Y g\nu$ for any $f\in L^1(\mu), g\in L^1(\nu)$. Then, for all $A\in\mathcal{B}(X)$ we have
    $$
    \pi(A\times Y) 
    = \int_{X\times Y} 1_{A\times Y}(x,y)d\pi(x,y) 
    = \int_{X\times Y} 1_A(x)d\pi(x,y)
    = \int_X 1_A \mu
    = \mu(A).
    $$
    Similarly, it holds $\pi(X\times B) = \nu(B)$ for all $B\in\mathcal{B}(Y)$. Then, $\pi\in Cpl(\mu, \nu)$ is a coupling between $\mu$ and $\nu$. \\
    Conversely, suppose that $\pi\in Cpl(\mu, \nu)$ is a coupling between $\mu$ and $\nu$, and let $f\in L^1(\mu)$ be an integrable function. Then, 
    \begin{align*}
        \int_{X\times Y} f(x) d\pi(x,y) 
        &= \int_{p^{-1}_X(X)} f(p_X(x,y)) d\pi(x,y) = \\
        &= \int_X f(z) d(p_X\pi)(x) = \text{ by substitution with } z = p_X(x,y) \\
        &= \int_X f(z) d\mu(z) \text{ since } \pi\in Cpl(\mu, \nu).
    \end{align*}
    Similarly, one can prove that $\int_{X\times Y} gd\pi = \int_Y gd\nu$ for all $g\in L^1(\nu)$.
\end{proof}

\begin{example}
    Let $\mu, \nu\sim\mathcal{N}(0,1)$ be univariate standard normal distributions. Then, any coupling of $\mu$ and $\nu$ is a bivariate normal distribution $$\pi\sim\mathcal{N}\left((0,0)^T, \begin{pmatrix}1 & \rho \\ \rho & 1\end{pmatrix}\right)$$ for some correlation parameter $\rho\in[-1, 1]$.
\end{example}

\begin{remark}
    Intuitively, a coupling of two measures $\mu, \nu$ encodes a specific dependency structure between the two distribution. In particular, $Cpl(\mu, \nu)\neq\varnothing$ for any couple of probability measures, because we can always consider the product distribution $\mu\otimes\nu\in Cpl(\mu, \nu)$ which encodes the independence structure. Note that, by definition of product measure, it holds
    $$
    (\mu\otimes\nu)(X\times B) = \mu(X)\nu(B)=\nu(B), \text{ and } \nu\otimes\nu(A\times Y) = \mu(A)\nu(Y)=\nu(A).
    $$
\end{remark}

\begin{example}
\label{example:dirac_coupling}
    Let $\delta_x, \ \delta_y$ be Dirac measures on $X, Y$ respectively. Then, there is a unique coupling between $\delta_x$ and $\delta_y$, given by the independent coupling $\delta_x \otimes \delta_y\in Cpl(\delta_x, \delta_y)$.\\
    Indeed, any coupling $\pi\in Cpl(\delta_x, \delta_y)$ satisfies 
    $$
    \pi((X\setminus\{x\})\times Y) = \delta_x(X\setminus\{x\}) = 0, 
    \quad
    \pi(X\times (Y\setminus\{y\})) = \delta_y(Y\setminus\{y\}) = 0
    $$
    Then, the support of $\pi$ is contained in $\{(x, y)\}$, which implies $\pi = \delta_{(x, y)} = \delta_x\otimes\delta_y$.
\end{example}

\begin{lemma}
\label{lemma:coupling_composition}
    Let $f:X\to Y, g: X\to Z$ be measurable maps, and $\mu\in\mathcal{P}(X)$ probability measure on $X$. Then,
    $$
    (f, g)_\#\mu \in Cpl(f_\#\mu, g_\#\mu),
    $$
    where $(f,g): X \to Y\times Z \text{ such that } (f,g)(x) = (f(x), g(x))$.
\end{lemma}
\begin{proof}
    Let $B\in\mathcal{B}(Y)$, and we prove that the distribution induced by projection of $(f, g)_\#\mu$ on $Y$ is the measure $f_\#\mu$. We have
    \begin{align*}
        p_{Y_\#}((f, g)_\#\mu)(B) 
        &= ((f, g)_\#\mu)(p^{-1}_Y(B)) = \\
        &= \mu((f,g)^{-1}(p^{-1}_Y(B))) = \\
        &= \mu(f^{-1}(B)) = (f_\#\mu)(B) \text{ since } p_Y \circ (f,g) = f.
    \end{align*}
    Similarly, it can be proven that the second marginal of $(f, g)_\#\mu$ is the measure $g_\#\mu$.
\end{proof}

\begin{proposition}
    Let $\mu, \nu$ be probability measures on $X, Y$ respectively. Then, the map $T\mapsto (Id_X, T)_\#\mu$ gives an injection $$\mathcal{T}(\mu, \nu) \hookrightarrow Cpl(\mu, \nu).$$
\end{proposition}
\begin{proof}
    Let $T: X\to Y$ be a transport map from $\mu$ to $\nu$, and consider the measurable map $$(Id_X, T) : X \to X\times Y \text{ such that } (Id_X, T)(x) := (x, T(x)).$$
    Then, the probability measure $$\pi := (Id_X, T)_\#\mu \in Cpl(Id_\#\mu, T_\#\mu) = Cpl(\mu, \nu)$$ is a coupling between $\mu$ and $\nu$ by (\ref{lemma:coupling_composition}).
\end{proof}

\begin{remark}
    The mass of the coupling $\pi=(Id_X, T)_\#\mu$ is concentrated on the graph of $T: X\to Y$. Equivalently, the coupling induced by transport maps are the laws of the random vectors of the type $(X, T(X))$. Hence, in the analogy where a transport maps moves entirely the mass of a point $x\in X$ to a point $y=T(x)\in Y$, we can interpret couplings as randomized transport maps. Indeed, a coupling $\pi\in Cpl(\mu, \nu)$ can be interpreted as a random map which assigns the map of the point $x\in X$ to a point $y\in Y$ with probability $\pi(Y\in \cdot \mid X=x)$. For this reason, we will sometimes refer to the couplings as transport plans.
\end{remark}

\begin{example} \hfill
    \begin{enumerate}[(i)]
        \item Let $\mu, \nu\in\mathcal{P}(\R)$ be probability measures with $\mu$ continuous. Then, the Monge transport $q_\nu \circ F_\mu : \R \to \R$ from $\mu$ to $\nu$ induces via canonical injection $\mathcal{T}(\mu, \nu)\hookrightarrow Cpl(\mu, \nu)$ a transport plan 
        $$
        \pi:=(Id_\R, q_\nu\circ F_\mu)_\# \mu \in Cpl(\mu, \nu),
        $$
        called monotone coupling. Intuitively, the monotone coupling between $\mu$ and $\nu$ assigns the masses in a monotone way, since both $q_\nu$ and $F_\mu$ are monotone. Moreover, the support of the coupling is a monotone subset of $\R^2$, since 
        $$
        (x_1 - x_2)(T(x_1) - T(x_2)) \geq 0
        $$
        for any couple $(x_1, T(x_1)), \ (x_2, T(x_2))$ in the support of $\pi$.
        \item Let $\mu, \nu\in\mathcal{P}(\R)$ be probability measures on the real line, and consider the probability measure $\pi:=(q_\mu, q_\mu)_\# \lambda$ on $\R^2$ induced by pushforward of the Lebesgue measure under the map 
        $$
        (q_\mu, q_\nu) : [0,1] \to \R^2 \text{ such that } (q_\mu, q_\nu)(u) = (q_\mu(u), q_\nu(u))
        $$
        Then, $\pi:=(q_\mu, q_\mu)_\# \lambda \in Cpl(\mu, \nu)$ is a coupling between $\mu$ and $\nu$ by Lemma \ref{lemma:coupling_composition}, called comonotone coupling.
    \end{enumerate}
\end{example}

\begin{remark}
    The previous examples suggest a connection between couplings and copulas. Indeed, copulas provide a parametrization of couplings, while optimal transport selects special couplings according to geometric or variational criteria.
\end{remark}



\subsection{Topological Properties of $\mathcal{T}(\mu, \nu), \ Cpl(\mu, \nu)$}
We have seen that the space of transport maps $\mathcal{T}(\mu, \nu)$ can be interpreted as a subspace of the space of couplings $Cpl(\mu, \nu)$ through the mapping $T\mapsto (Id_X, T)_\#\mu$. In the next section, we will define the Kantorovich and the Monge problem as optimization problems on the spaces $\mathcal{T}(\mu, \nu), \ Cpl(\mu, \nu)$ respectively. Then, we will now investigate some properties of these spaces. 

\begin{definition}
    A sequence of measures $\{\mu_n\}_n\subseteq \mathcal{P}(X)$ converges weakly to a measure $\mu\in\mathcal{P}(X)$ if $$\int_X fd\mu_n \underset{n\to\infty}{\longrightarrow} \int_X f d\mu$$ for any continuous bounded function $f:X\to\R$, i.e. $f\in C_b(X)$ .
\end{definition}

The notion of weak convergence induces a topology on the space of probability measures $\mathcal{P}(x)$, called weak topology, where the closed subspaces $C\subseteq\mathcal{P}(X)$ are exactly the ones preserving the limits.

\begin{definition}
    A family of probability measures $\mathcal{A}\subseteq\mathcal{P}(X)$ is relatively compact respect to the weak topology if the closure $\bar{\mathcal{A}}\subseteq\mathcal{P}(X)$ is compact. Equivalently, $\mathcal{A}$ is relatively compact if each sequence $\{\mu_n\}_n\subseteq\mathcal{A}$ admits a convergent subsequence in $\mathcal{P}(X)$.
\end{definition}

\begin{definition}
    A family of probability measures $\mathcal{A}\subseteq\mathcal{P}(X)$ is tight if for all $\epsilon>0$ there exists a compact subspace $K_\varepsilon\subseteq\mathcal{P}(X)$ such that $$\sup_{\mu\in\mathcal{A}} \mu(X\setminus K_\varepsilon)<\varepsilon.$$
\end{definition}

The next result gives a practical characterization of the notion of relative compacteness in $\mathcal{P}(X)$, proving that it is indeed equivalent to the tightness.

\begin{theorem}
\label{thm:prokhorov}
    (\textbf{Prokhorov}) Let $\mathcal{A}\subseteq\mathcal{P}(X)$ be a family of probability measures, then $$\mathcal{A}\subseteq\mathcal{P}(X) \text{ relatively compact resp. weak topology } \iff \mathcal{A}\subseteq\mathcal{P}(X) \text{ tight.}$$
\end{theorem}

\begin{lemma}
\label{lemma:tight_cartesian_product}
    $\mathcal{A}_1\subseteq\mathcal{P}(X), \mathcal{A}_2\subseteq\mathcal{P}(Y)$ tight families of probability measures. Then, the family $$\mathcal{A}:= \{\pi\in\mathcal{P}(X\times Y) \mid p_X\pi\in\mathcal{A}_1, \ p_Y\pi\in\mathcal{A}_2\}$$ is tight.
\end{lemma}
\begin{proof}
    Let $\pi\in\mathcal{A}$ with marginals $\mu:=p_X\pi\in\mathcal{A}_1, \ \nu:=p_Y\pi\in\mathcal{A}_2$, and consider the compact sets $K_1\subseteq X, \ K_2\subseteq Y$ such that $\mu(X\setminus K_1)<\varepsilon$, $\nu(Y \setminus K_2)<\varepsilon$ for a fixed $\varepsilon>0$. Then, for the compact $K_1 \times K_2\subseteq X \times Y$, we have
    \begin{align*}
        \pi((X\times Y) \setminus (K_1 \times K_2))
        &\leq \pi(((X\setminus K_1)\times Y) \cup (X\times (Y\setminus K_2)))\leq \\
        &\leq \pi((X\setminus K_1)\times Y) + \pi(X\times (Y\setminus K_2)) \leq \\
        &\leq \mu(X\setminus K_1) + \nu(Y\setminus K_2)< 2\varepsilon \quad \text{ since } \pi\in Cpl(\mu, \nu) \\
    \end{align*}
    Then, the thesis follows by taking LHS supremum.
\end{proof}

\begin{lemma}
\label{lemma:integral_linear}
    Let $\mu, \nu\in\mathcal{P}(X)$ be probability measures on $X$ such that $a\mu + b\nu\in\mathcal{P}(X)$ for some $a,b\geq 0$. If $f\in L^1(\mu)\cap L^1(\nu)$, then $f\in L^1(a\mu+b\nu)$ with
    $$
    \int_X f \ d(a\mu+b\nu) = \int_X f \ d\mu + \int_X f \ d\nu. 
    $$
    In particular, the map $\mu \mapsto \int_X f \ d \mu$ is linear with respect to convex combinations of measures.
\end{lemma}
\begin{proof}
    Firstly, suppose that $f = \sum_{i=1}^n c_i 1_{A_i}$ simple function for some $c_i\geq 0$, and $A_i\in\mathcal{B}(X)$ for $i=1, \dots, n$. Then, we get 
    \begin{align*}
        \int f \ d(a\mu + b\nu) 
        = \sum_{i=1}^n c_i (a \mu(A_i) + b\nu(A_i)) 
        = a \int f \ d\mu + b \int f \ d\nu
    \end{align*}
    Now, suppose that $f = \lim_{n\to \infty} f_n$ with $\{f_n\}_n$ increasing sequence of simple functions. Then, by monotone convergence we obtain
    $$
    \int f \ d(a\mu + b\nu) 
    = \lim_{n\to\infty} \int f_n \ d(a\mu+b\nu)
    = a\int f \ d\mu + b\int f \ d\nu
    $$
    Finally, we get the thesis by decomposing $f=f^+ - f^-$ and applying the previous step to $f^+, f^-$.
\end{proof}

The next result states the properties of the space $Cpl(\mu, \nu)$ of couplings between $\mu$ and $\nu$. 

\begin{proposition}
\label{prop:Cpl()_propr}
    Let $\mu, \nu$ be probability measures on $X, Y$ respectively. Then, the following holds:
    \begin{enumerate}[(i)]
        \item $Cpl(\mu, \nu)$ convex.
        \item $Cpl(\mu, \nu)\subseteq\mathcal{P}(X\times Y)$ compact.
    \end{enumerate}
\end{proposition}
\begin{proof}
    \textit{(i)} Let $\pi_1, \pi_2\in Cpl(\mu, \nu)$ be couplings between $\mu$ and $\nu$, and we prove that $\alpha\pi_1 + (1-\alpha)\pi_2$ is a coupling between $\mu$ and $\nu$. Let $f\in L^1(\mu)$ be an integrable function, then
    \begin{align*}
        \int_{X\times Y} &f(x) d(\alpha\pi_1 + (1-\alpha)\pi_2)(x, y) = \\
        &= \alpha\int_{X\times Y} f(x)d\pi_1(x,y) \ + \ (1-\alpha) \int_{X\times Y} f(x)d\nu(x,y) = \\
        &= \alpha\int_X f(x)d\mu(x) \ + \ (1-\alpha)\int_X f(x)d\pi_2(x, y) = \text{ by (\ref{lemma:integral_linear}) and  (\ref{prop:characterization_cpl})} \\
        &= \int_X f(x)d\mu(x).
    \end{align*}
    Analogously, the equality  holds for the other marginal $\nu$, proving that the probability measure $\alpha\pi_1 + (1-\alpha)\pi_2 \in Cpl(\mu,  \nu)$ is a coupling. \\

    \noindent
    \textit{(ii)} Consider the compact families $\mathcal{A}_1 = \{\mu\}\subseteq\mathcal{P}(X)$ and $\mathcal{A}_2 = \{\nu\}\subseteq\mathcal{P}(Y)$, which are tight by Prokhorov's theorem. Then, the space of coupling between $\mu$ and $\nu$ $$Cpl(\mu, \nu) =\{\pi\in\mathcal{P}(X\times Y) \mid p_X\pi=\mu, \ p_Y\pi=\nu\}$$ is tight by Lemma \ref{lemma:tight_cartesian_product}.
    In particular, $Cpl(\mu, \nu)\subseteq\mathcal{P}(X\times Y)$ is relatively compact respect to the weak topology, by Theorem \ref{thm:prokhorov}.\\
    Now, consider a sequence of probability measures $\{\pi_n\}_n\subseteq Cpl(\mu, \nu)$ converging weakly to a probability measure $\pi\in\mathcal{P}(X\times Y)$. We prove that $Cpl(\mu, \nu)$ is closed in $\mathcal{P}(X\times Y)$, i.e. we prove that the limit measure $\pi$ is a coupling between $\mu$ and $\nu$. \\
    Let $f \in C_b(X)$ a continuous and bounded function, then we have
    \begin{align*}
        \int_{X\times Y} f(x) d\pi(x,y) &= \int_{X\times Y} \tilde f(x, y)d\pi(x,y) = \text{ with } \tilde f(x,y):=f(x) \\
        &= \lim_{n\to\infty} \int_{X\times Y} \tilde f(x,y)d\pi_n(x,y) = \text{ since } \tilde f\in C_b(X\times Y) \\
        &= \lim_{n\to\infty} \int_{X\times Y} f(x)d\mu(x) 
        = \int_X f(x)d\mu(x) \text{ by Proposition \ref{prop:characterization_cpl}.} 
    \end{align*}
    Similarly, one can prove that $\int_{X\times Y} g(x) d\pi(x,y) = \int_Y g(x)d\nu(x)$ for all continuous and bounded $g\in C_b(Y)$. Then, $\pi$ is a coupling between $\mu$ and $\nu$ by Proposition \ref{prop:characterization_cpl}, and $Cpl(\mu, \nu)$ is closed. Hence, $Cpl(\mu, \nu)$ is also compact as a closed and relatively compact subspace of $\mathcal{P}(X\times Y)$.
\end{proof}

\begin{remark}
    In general, the space $\mathcal{T}(\mu, \nu)$ of transport maps from $\mu$ to $\nu$ does not satisfy any of the properties above.
    \begin{enumerate}[(i)]
        \item \textit{$\mathcal{T}(\mu, \nu)$ not convex.} \\
        Consider the measures $\mu := \frac{1}{2}(\delta_{-1} + \delta_1), \ \nu := \frac{1}{2}(\delta_{-2} + \delta_{2}) \in\mathcal{P}(\R)$, and the measurable maps $T_1, T_2 : \R\to\R$ such that $T_1(x) = 2x, \ T_2(x)=-2x$. Then, we have 
        $$
        T_1\mu = \nu,  \quad T_2\mu = \nu
        $$
        since $T_1(-1)=-2, \ T_1(1)=2$ and $T_2(-1)=2, \ T_2(1)=-2$. However, the convex combination $T:= \frac{1}{2}(T_1 + T_2)$ is not a transport from $\mu$ to $\nu$, since $T(-1)=T(1)=0$.    
        
        \item \textit{$\mathcal{T}(\mu, \nu)$ not closed respect to the weak topology.} \\
        Let $\mu=\nu=\lambda$ be the Lebesgue measure on the interval $[0,1]$, and consider the sequence of measurable maps 
        \begin{align*}
            T_n: [0,1]\to[0,1] \
            \text{ such that } \
            T_n(x) = 
            \begin{cases}
                x & x\in\bigcup_{k\in\N} [\frac{2k}{2n}, \frac{2k+1}{2n}]\\
                1-x & x\in\bigcup_{k\in\N} [\frac{2k+1}{2n}, \frac{2k+2}{2n}]
            \end{cases}
        \end{align*}
        Then, each $T_n\in\mathcal{T}(\mu, \nu)$ is a transport maps from the Lebesgue measure to itself. Moreover, note that the induced coupling $\pi_n:=(Id_{[0,1]}, T_n)_\#\mu$ converges weakly to 
        $$
        \pi:=\frac{1}{2}(Id_{[0,1]}, Id_{[0,1]})_\#\mu + (Id_{[0,1]}, 1- Id_{[0,1]})_\#\mu)
        $$
        because half of the mass is sent to $y=x$ and the other half is sent to $y=1-x$. However, $\pi$ is not induced by any transport map, since its support in $[0,1]^2$ is not the graph of a function. Therefore $\pi$ is a proper coupling between $\mu$ and $\nu$, and $\pi\notin\mathcal{T}(\mu, \nu)$. 
    \end{enumerate}
\end{remark}

\subsection{Kantorovich and Monge Problems}
In this section, we finally the optimal transport problems, both in the Monge formulation and the Kantorovich formulation, and then we will study some properties of these problems.\\
The first notion we introduce models the cost of transporting one unity of mass from a point $x$ of the domain space $X$ to another point $y$ in the arrival space $Y$.

\begin{definition}
    A cost function for an optimal transport problem with distributions $\mu\in\mathcal{P}(X), \ \nu\in\mathcal{P}(Y)$ is a measurable map $c:X\times Y\to\extRplus$. \\
    In particular, we say that $c$ is bounded below if $$c(x, y)\geq a(x) + b(y) \text{ for all } x\in X, y\in Y, $$ for some upper-semicontinuous functions $a: X\to \extRminus, \ b: Y\to \extRminus$ with $a\in L^1(\mu), \ b\in L^1(\nu)$.
\end{definition}

In many practical examples, we will consider the setting $X=Y$ with cost given by the metric $d: X\times X\to [0, \infty)$ induced on the Polish space $X$.

\begin{remark}
    The lower bound assumption $c(x, y)\geq a(x) + b(y)$ guarantees that the expected cost $\E[c(X, Y)]$ is well defined in $\extRplus$, with $X\sim \mu, \ Y \sim \nu$. \\
    In most of the cases, we will consider non-negative cost functions $c: X\times Y\to[0, +\infty]$, i.e. with $a=0, b=0$.
\end{remark}

\begin{definition}
    The Monge functional with cost $c:X\times Y\to\extRplus$ is the map $P^M_c(\mu, \nu): \mathcal{P}(X)\times\mathcal{P}(Y)\to\extRplus$ such that
    $$
    P^M_c(\mu, \nu) := \inf_{T\in\mathcal{T}(\mu, \nu)} \int_X c(x, T(x)) \ d\mu(x) 
    \quad (MP)_c
    $$
    In particular, a transport map $T\in\mathcal{T}(\mu, \nu)$ is called optimal transport if it attains the infimum of (MP).
\end{definition}

\begin{definition}
    The Kantorovich functional with cost $c:X\times Y\to\extRplus$ is the map $P^K_c(\mu, \nu): \mathcal{P}(X)\times\mathcal{P}(Y)\to\extRplus$ such that
    $$
    P^K_c(\mu, \nu) := \inf_{\pi\in Cpl(\mu, \nu)} \int_{X\times Y} c \ d\pi 
    \quad (KP)_c
    $$
    In particular, a coupling $\pi\in Cpl(\mu, \nu)$ is called optimal transport plan if it attains the infimum of $(KP)_c$.
\end{definition}

Intuitively, both problems represent a minimization problem of an expected cost, depending on some marginal distributions $\mu, \nu$. Indeed, we can rewrite the functionals $P^M_c(\mu, \nu)$ and $P^K_c(\mu, \nu)$ as
$$
P^M_c(\mu, \nu) := \inf_{\substack{T\in\mathcal{T}(\mu, \nu), \\ X\sim\mu}} \E[c(X, T(X))], 
\quad 
P^K_c(\mu, \nu) := \inf_{\substack{\pi\in Cpl(\mu, \nu) \\ X\sim\mu, Y\sim\nu}} \E_\pi[c(X, Y)].
$$

\begin{remark}
    The Kantorovich problem is a generalization of the Monge problem, since each transport map can be interpreted as a coupling of measures via canonical injection. However, (KP) is much better behaved than (MP), as we saw in the previous section:
    \begin{itemize}
        \item $Cpl(\mu, \nu)$ is non-empty, convex, and compact resp. weak topology.
        \item $\mathcal{T}(\mu, \nu)$ is in general not convex and not closed resp.\ weak topology
    \end{itemize}
    In particular, the convexity of (KP) endows the problem with a richer structure, allowing the use of a wide range of tools from convex analysis.
\end{remark}

\subsection{Existence}
The first natural question arising from the definitions of (KP) and (MP) is the existence of optimizers, which will be clarified in the next theorem.

\begin{definition}
    A function $f:X\to \bar\R$ is lower-semicontinuous (LSC) if $$\liminf_{n\to\infty} f(x_n)\geq f(x)$$ for all sequences $\{x_n\}_n\subseteq X$ converging to $x\in X$.
\end{definition}

\begin{lemma}
\label{lemma:lsc_propr}
    Let $f: X\to\bar\R$ be a function, then the following hold:
    \begin{enumerate}[(i)]
        \item $f \ LSC \iff \{f\leq y\}\subseteq X \text{ closed } \iff \{f > y\}\subseteq X \text{ open}$.
        \item If $f \ LSC $, then $f$ attains minimum on any $K\subseteq X$ compact.
        \item (Baire Theorem) $f \ LSC \iff f =\lim_{n\to\infty} f_n \text{ for some } \{f_n\}_n\subseteq C(X), \ f_n\uparrow f$.
    \end{enumerate}
\end{lemma}

\begin{lemma}
\label{lemma:integral_cost_lsc}
    Let $c: X\times Y\to \extRplus$ be a lower-semicontinuous cost function such that $c\geq h$ for some upper-semicontinuous function $h:X\times Y\to \extRminus$. If $\{\pi_n\}_n$ sequence of probability measures in $\mathcal{P}(X\times Y)$ converging weakly to a probability $\pi$ such that $h\in L^1(\pi_n), \ h\in L^1(\pi)$, and $$\int_{X\times Y} h \ d\pi_n\underset{n\to\infty}{\longrightarrow} \int_{X\times Y} h \ d\pi,$$ then $$\int_{X\times Y} c \ d\pi \leq \liminf_{n\to\infty} \int_{X\times Y} c \ d\pi_n.$$ 
    In particular, if $c\geq 0$ then the map $$\pi\mapsto\int_{X \times Y} c \ d\pi$$ is lower-semicontinuous.
\end{lemma}
\begin{proof}
Note that we may assume that $c$ is non-negative, otherwise we can replace $c$ with $c - h\geq 0$. Then, $c$ can be written
as the pointwise limit of a non decreasing family $\{c_k\}_k$ of continuous and bounded functions by Lemma \ref{lemma:lsc_propr}, i.e. $c_k \uparrow c$. Then, we have
\begin{align*}
    \int_{X\times Y} c \ d\pi 
    &= \lim_{k\to\infty} \int_{X\times Y} c_k \ d\pi = \text{ by monotone convergence} \\
    &= \lim_{k\to\infty} \lim_{n\to\infty}\int_{X\times Y} c_k \ d\pi_n \text{ by weak convergence}
\end{align*}
Now, note that 
$$
\int_{X\times Y} c_k \ d\pi_n \leq \int_{X\times Y} c \ d\pi_n \text{ for all } n\in\N
$$
since $\{c_n\}_n$ non decreasing sequence converging to $c$. Hence, by taking the liminf for $n\to \infty$ we get for all $k\in\N$
$$
\int_{X\times Y} c_k \ d\pi 
= \lim_{n\to\infty} \int_{X\times Y} c_k \ d\pi_n
= \liminf_{n\to\infty} \int_{X\times Y} c_k \ d\pi_n 
\leq \liminf_{n\to\infty} \int_{X\times Y} c \ d\pi_n,
$$
by weak convergence of $\pi_n$ to $\pi$. Finally, the thesis follows by taking the lim for $k\in\N$
$$
\lim_{k\to\infty}\lim_{n\to\infty} \int_{X\times Y} c_k \ d\pi_n 
\leq \liminf_{n\to\infty} \int_{X\times Y} c \ d\pi_n.
$$
\end{proof}

\begin{theorem}
\label{thm:optimal_plan_existence}
    (Optimal Transport Plans Existence) Let $c: X\times Y\to \extRplus$ be a lower-semicontinuous and bounded below cost function for a Kantorovich problem with marginals $\mu\in\mathcal{P}(X), \ \nu\in\mathcal{P}(Y)$. Then, there exists an optimal transport plan for the Kantorovich problem $(KP)_c$, i.e.
    $$
    \exists \ \pi^*\in\argmin_{\pi\in Cpl(\mu, \nu)} \int_{X\times Y} c \ d\pi.
    $$
\end{theorem}
\begin{proof}
    Let $\{\pi_n\}_n\subseteq Cpl(\mu, \nu)$ be a sequence of couplings between $\mu$ and $\nu$ with expected cost converging to the infimum average cost, i.e. such that
    $$
    \int_{X\times Y} c \ d\pi_n \underset{n\to\infty}{\longrightarrow} P^K_c(\mu, \nu) = \inf_{\pi\in Cpl(\mu, \nu)} \int_{X\times Y} c \ d\pi 
    $$
    Then, there exists a subsequence $\{\pi_{n_k}\}_k \subseteq Cpl(\mu, \nu)$ converging weakly to a coupling $\pi\in Cpl(\mu, \nu)$, since $Cpl(\mu, \nu)$ is compact by Proposition \ref{prop:Cpl()_propr}. Moreover, note that $c\geq h$ with $h(x, y) = a(x) + b(y)$, and it also holds $h\in L^1(\pi_n), \ h\in L^1(\pi)$ with
    $$
    \int_{X\times Y} h \ d\pi_n 
    = \int_{X\times Y} a\ d\mu + \int_{X\times Y} b\ d\nu 
    = \int_{X\times Y} c\ d\pi 
    $$
    by Proposition \ref{prop:characterization_cpl}, since $\pi_n, \pi\in Cpl(\mu, \nu)$. \\
    Then, we get that $\pi$ attains the infimum by Lemma \ref{lemma:integral_cost_lsc} since the cost is LSC, i.e.
    $$
    \int_{X\times Y} c \ d\pi \leq \liminf_{n\to\infty} \int_{X\times Y} c \ d\pi_n = P^K_c(\mu, \nu).
    $$
\end{proof}

\begin{remark}
    The existence theorem does not imply that the optimal cost is finite. It might be that all transport plans lead to an infinite total cost, i.e. $\int_{X\times Y} c \ d\pi = +\infty$ for all $\pi\in Cpl(\mu, \nu)$. A simple condition to rule out this annoying possibility is
    $$
    \int_{X\times Y} c(x, y) \ d\mu(x) \ d\nu(y) < +\infty,
    $$ 
    which guarantees that at least the independent coupling $\mu\otimes\nu$ has finite total cost. \\
    Alternatively, one can make the stronger assumption
    $$
    c(x, y) \leq c_X(x) + c_Y(y) \text{ for some functions } c_X\in L^1(\mu), \ c_Y\in L^1(\nu). 
    $$
\end{remark}

\subsection{Duality}
The previous result guarantees the existence of optimizer for the Kantorovich problem with rather mild assumptions on the regularity of the cost function $c$. Once existence is established, some other natural questions arise:
\begin{itemize}
    \item What is the dual formulation of the convex problem (KP)?
    \item Is the solution of (KP) unique?
    \item Do (MP) and (KP) have the same solutions?
    \item Is there any structure in the optimizer $\pi^*$ of (KP)?
\end{itemize}

\begin{definition}
    The dual Kantorovich functional with cost $c:X\times Y\to\extRplus$ is the map $D^K_c(\mu, \nu): \mathcal{P}(X)\times\mathcal{P}(Y)\to\extRplus$ such that
    $$
    D^K_c(\mu, \nu) := \sup_{(\varphi, \psi)\in D(c)} \left(\int_X \varphi \ d\mu \ + \ \int_Y \psi \ d\nu \right) 
    \quad \text{(DP)}
    $$
    with $D(c) := \{(\varphi, \psi)\in L^1(\mu)\times L^1(\nu) \mid  c \geq \varphi \oplus \psi\}$ research space.
\end{definition}

Intuitively, if the Kantorovich problem (KP) aims to minimize some average cost, then the dual problem (DP) aims to maximize the some profit.\\
In the following we will use two distinct notions of duality, and we will say that:
\begin{itemize}
    \item Weak Duality holds if $D^K_c(\mu, \nu) \leq P^K_c(\mu, \nu)$
    \item Strong Duality holds if $D^K_c(\mu, \nu) = P^K_c(\mu, \nu)$
\end{itemize}

\begin{proposition}
\label{prop:weak_duality}
    (Weak Duality) Let $c: X\times Y\to \extRplus$ be a cost function for a Kantorovich problem with marginals $\mu\in\mathcal{P}(X), \ \nu\in\mathcal{P}(Y)$. Then, weak duality holds, i.e. $D^K_c(\mu, \nu) \leq P^K_c(\mu, \nu)$.
\end{proposition}
\begin{proof}
    Let $\varphi\in L^1(\mu), \psi\in L^1(\nu)$ be integrable functions such that $\varphi\oplus \psi \leq c$, and consider a coupling $\pi\in Cpl(\mu, \nu)$. Then, we have
    \begin{align*}
        \int_X \varphi \ d\mu \ + \ \int_Y \psi \ d\nu 
        &= \int_{X\times Y} (\varphi \oplus \psi) \ d\pi \text{ by Proposition \ref{prop:characterization_cpl}} \\
        &\leq \int_{X\times Y} c \ d\pi.
    \end{align*}
    Hence, the thesis follows by taking the infimum of RHS over $\pi\in Cpl(\mu, \nu)$, and the supremum of the LHS over $(\varphi, \psi)\in D(c)$, that is
    $$
    D^K_c(\mu, \nu) = \sup_{(\varphi, \psi)\in D(c)} \left(\int_X \varphi \ d\mu \ + \ \int_Y \psi \ d\nu \right) 
    \leq 
    \inf_{\pi\in Cpl(\mu, \nu)} \int_{X\times Y} c \ d\pi = P^K_c(\mu, \nu)
    $$
\end{proof}

\begin{remark}
    Strong duality may not hold. \\ 
    For example, consider $X=Y=[0,1]$ with (non lower-semicontinuous) cost
    \begin{align*}
        c(x, y) = 
        \begin{cases}
            0 & x<y \\
            1 & x=y \\
            +\infty & \text{otherwise}.
        \end{cases}
    \end{align*}
    Then, duality does not hold for the Kantorovich problem with marginals $\mu=\nu=\lambda_{|_{[0,1]}}$. 
    Indeed, the identity coupling $\pi:=(Id_{[0,1]}, Id_{[0,1]})\lambda \in Cpl(\lambda, \lambda)$ realizes the infimum of the Kantorovich problem
    $$
    \inf_{\pi\in Cpl(\lambda, \lambda)} \int_{[0,1]^2} c \ d\pi = 1,
    $$
    since $\pi$ supported on the diagonal of $[0,1]^2$. However, the dual problem satisfies
    $$
    \sup_{\varphi\oplus\psi \leq c} \int_{[0,1]} \varphi(x)+\psi(x) \ dx = 0.
    $$
\end{remark}

In order to establish a theorem for the Strong Duality, we first need the following technical result.

\begin{lemma}
\label{lemma:min_max_principle}
    (Min-Max Principle) Let $H_1, H_2$ be vector spaces with $H_1$ locally compact, and let $K\subseteq H_1$ be a convex and compact subspace, $L\subseteq H_2$ be a convex subspace. Let $F: K\times L\to \R$ be a map such that
    \begin{enumerate}[(i)]
        \item $x\mapsto F(x, y)$ continuous and convex for all $y\in L$
        \item $y\mapsto F(x, y)$ concave for all $x\in K$
    \end{enumerate}
    Then, 
    $$
    \sup_{y\in L}\inf_{x\in K} F(x, y) = \inf_{x\in K}\sup_{y\in L} F(x, y)
    $$
\end{lemma}
\begin{proof}
    Theorem 37.1, \cite{rockafellar1997convex}.
\end{proof}

\begin{theorem}
\label{thm:strong_duality} (Duality, I)
    Let $c: X\times Y\to \extRplus$ be a lower-semicontinuous and bounded below cost function for a Kantorovich problem with marginals $\mu\in\mathcal{P}(X), \ \nu\in\mathcal{P}(Y)$. Then, the following hold: 
    \begin{enumerate}[(i)]
        \item Strong duality holds, i.e. $$D^K_c(\mu, \nu) = P^K_c(\mu, \nu).$$
        \item The research space $D(c)$ of the dual problem can be restricted to $$\tilde D(c) := \{(\varphi, \psi)\in C_b(X)\times C_b(Y)\mid c\geq \varphi\oplus\psi\},$$ i.e., it holds
    $$
    \sup_{\substack{(\varphi, \psi)\in L^1(\mu)\times L^1(\nu) \\ c\geq \varphi\oplus\psi}} \left(\int_X \varphi \ d\mu \ + \ \int_Y \psi \ d\nu \right)
    =
    \sup_{\substack{(\varphi, \psi)\in C_b(X)\times C_b(Y) \\ c\geq \varphi\oplus\psi}} \left(\int_X \varphi \ d\mu \ + \ \int_Y \psi \ d\nu \right).
    $$
    \item Let $\varphi^*\in L^1(\mu), \psi^*\in L^1(\nu)$ be optimizer of the dual problem, then
    $$
    \pi\in Cpl(\mu, \nu) \text{ optimal} \iff c = \varphi^* \oplus \psi^* \quad \pi-a.s.
    $$
    Similarly, if $\pi^*$ is an optimizer of the Kantorovich problem, then
    $$
    \varphi\in L^1(\mu), \ \psi\in L^1(\nu) \text{ s.t. $\varphi\oplus\psi\leq c$ optimal } \iff c = \varphi \oplus \psi \quad \pi^*-a.s.
    $$
    \end{enumerate}
\end{theorem}
\begin{proof} \ \\
    \textit{(i)} We will prove the result in the particular case where $X, Y$ are compact spaces and the cost $c: X\times Y\to \extRplus$ is a continuous function. \\
    Let $\mu, \nu$ be probability measures on $X, Y$ respectively, and we prove that $P^K_c(\mu, \nu) = D^K_c(\mu, \nu)$. Consider the map $\chi : \mathcal{P}(X\times Y) \to \{0, \infty\}$ such that
    $$
    \chi(\pi) :=
    \begin{cases}
        0 & \pi\in Cpl(\mu, \nu) \\
        +\infty & \text{otherwise}
    \end{cases}
    $$
    Note that for a probability measure $\pi\in\mathcal{P}(X\times Y)$ we have by Proposition \ref{prop:characterization_cpl}
    $$
    \pi\in Cpl(\mu, \nu) \iff \int_{X\times Y} \varphi\oplus\psi \ d\pi = \int_X \varphi \ d\mu + \int_Y \psi \ d\nu.
    $$
    Then, we can rewrite the map $\pi\to\chi(\pi)$ as 
    \begin{align*}
        \chi(\pi) = \sup_{(\varphi, \psi)\in L^1(\mu)\times L^1(\nu)} 
        \left(\int_X \varphi \ d\mu + \int_Y \psi \ d\nu - \int_{X\times Y} \varphi \oplus \psi \ d\pi\right).
    \end{align*}
    Therefore, that Kantorovich problem is given by
    \begin{align*}
        P^K_c(\mu, \nu) 
        &= \inf_{\pi\in Cpl(\mu, \nu)} \int_{X\times Y} c \ d\pi
        = \inf_{\pi\in \mathcal{P}(X\times Y)} \left( \chi(\pi) + \int_{X\times Y} c \ d\pi \right) = \\
        &= \inf_{\pi\in\mathcal{P}(X\times Y)} \sup_{(\varphi, \psi)\in L^1(\mu)\times L^1(\nu)} \left(\int_X \varphi \ d\mu + \int_Y \psi \ d\nu + \int_{X\times Y} (c-\varphi \oplus \psi) \ d\pi\right)
    \end{align*}
    Now, note that the space $\mathcal{P}(X\times Y)$ is compact for $X, Y$ compact spaces, and we can apply the min-max principle with $K=\mathcal{P}(X\times Y)$ and $L=L^1(\mu)\times L^1(\nu)$. Thus, we have
    \begin{align*}
        P^K_c(\mu, \nu) 
        &= \inf_{\pi\in\mathcal{P}(X\times Y)} \sup_{(\varphi, \psi)\in L^1(\mu)\times L^1(\nu)} \left(\int_X \varphi \ d\mu + \int_Y \psi \ d\nu + \int_{X\times Y} (c-\varphi \oplus \psi) \ d\pi\right) \\
        &= \sup_{(\varphi, \psi)\in L^1(\mu)\times L^1(\nu)} \inf_{\pi\in\mathcal{P}(X\times Y)} 
        \left(\int_X \varphi \ d\mu + \int_Y \psi \ d\nu + \int_{X\times Y} (c-\varphi \oplus \psi) \ d\pi\right) \\
        &\leq \sup_{(\varphi, \psi)\in L^1(\mu)\times L^1(\nu)} \inf_{(x, y)\in X\times Y}
        \left(\int_X \varphi \ d\mu + \int_Y \psi \ d\nu + c(x, y) - \varphi(x) - \psi(y)\right) \\
        &= \sup_{\substack{(\varphi, \psi)\in L^1(\mu)\times L^1(\nu) \\ c\geq \varphi\oplus\psi}}
        \left(\int_X \varphi \ d\mu + \int_Y \psi \ d\nu\right) = D^K_c(\mu, \nu).
    \end{align*}
    More generally, when $X, Y$ are not compact and the cost $c$ is not continuous, we can approximate the cost function $c$ with a non-decreasing sequence of continuous functions $\{c_n\}_n$, and we can approximate $\mu, \nu$ with compactly supported probability measures $\tilde\mu, \tilde\nu$, and use the previous step. \\
    \textit{(ii)} Note that by density of $C_b(X)\subseteq L^1(\mu)$ and $C_b(Y)\subseteq L^1(\nu)$, the research space of maximizers can be restricted to the space $$\tilde D(c) := \{(\varphi, \psi)\in C_b(X)\times C_b(Y)\mid c\geq \varphi\oplus\psi\}.$$
    \textit{(iii)} Let $\varphi^*\in L^1(\mu), \ \psi^*\in L^1(\nu)$ be optimizer of the dual problem, and assume that $\pi\in Cpl(\mu, \nu)$ satisfies $c = \varphi^* \oplus \psi^*$ $\pi$-a.s. Then, we have 
    \begin{align*}
        \int_{X\times Y} c \ d\pi 
        &= \int_{X\times Y} \varphi^*\oplus\psi^* \ d \pi
        = \int_X \varphi^* \ d\mu + \int_Y \psi^* \ d\nu \\
        &= \sup_{\substack{(\varphi, \psi)\in L^1(\mu)\times L^1(\nu) \\ c\geq \varphi\oplus\psi}} \left(\int_X \varphi \ d\mu \ + \ \int_Y \psi \ d\nu \right) \\
        &= \inf_{\pi\in Cpl(\mu, \nu)} \int_{X\times Y} c \ d\pi \quad \text{by duality.}
    \end{align*}
    Hence, $\pi$ attains the infimum of the Kantorovich problem, and is then an optimizer.  
\end{proof}

\vspace{10pt}
\noindent
The next relevant result originates from a simple question: 
\begin{quote}
Given $(\varphi, \psi)\in L^1(\mu)\times L^1(\nu)$ such that $\varphi\oplus\psi\leq c$, achieving a certain score $\int_X \varphi \ d\mu + \int_X \psi \ d\nu$ in the dual problem, can we construct a couple $(\tilde \varphi, \tilde \psi)$ still satisfying $\tilde\varphi\oplus\tilde\psi\leq c$ and achieving a better core, i.e. such that 
$$
\int \varphi \ d\mu + \int \psi \ d\nu
\leq 
\int \tilde\varphi \ d\mu + \int \tilde\psi \ d\nu ?
$$
\end{quote}
For example, one may fix the function $\varphi$, and improve $\psi$ by considering the function
$$
\tilde\psi(y) := \inf_{x\in X} \ (c(x, y) - \varphi(x)).
$$
Then, we have
$$
\varphi(x) + \tilde \psi(y) 
= \varphi(x) + \inf_{\bar x\in X} \ (c(\bar x, y) - \varphi(\bar x))
\leq \varphi(x) + c(x, y) - \varphi(x) 
= c(x, y)
$$
and
$$
\int \varphi \ d\mu + \int \psi \ d\nu \leq \int\varphi \ d\nu + \int \tilde\psi \ d\nu 
$$
since $\psi\leq\tilde\psi$ by $\varphi \oplus \psi \leq c$.

\begin{definition}
    The transform of a measurable function $\varphi : X \to \R$ respect to a cost function $c: X \times Y \to \extRplus$, or $c$-transform of $\varphi$, is the map $\varphi^c: Y \to \R$ such that $$\varphi^c(y) := \inf_{x\in X} (c(x, y) - \varphi(x)).$$
    Similarly, the $c$-transform of a map $\psi: Y \to \R$ is the function $\psi^c : X \to \R$ such that $$\psi^c(x) := \inf_{y\in Y} (c(x, y) - \psi(y)).$$
    Moreover, we say that $\varphi$ is $c$-concave is $\varphi = \psi^c$ for some map $\psi$, and similarly $\psi$ is $c$-concave if $\psi = \varphi^c$ for some function $\varphi$.
\end{definition}

\begin{example} \hfill
    \begin{enumerate}[(i)]
        \item \textit{Quadratic Cost $c(x, y) = -x \cdot y$} \\
        Let $X= Y = \R^d$ with cost $c(x, y) = - x\cdot y$, and let $\varphi: \R^d \to \R$ be a measurable function. Then, 
        $$
        \varphi: \R^d \to \R \ c\text{-concave} \iff \varphi \ \text{concave and USC}
        $$
        In particular, the $c$-transform is given by
        $$
        \varphi^c(y) 
        = \inf_{x\in\R^d} \ (-x\cdot y - \varphi(x)) 
        = - \sup_{x\in\R^d} \ (x\cdot y - (-\varphi)(x))
        = - (-\varphi)^*(y),
        $$
        where $\varphi^*: \R^d \to \R$ Legendre transform of $\varphi$. \\
        Indeed, suppose that $\varphi= \psi^c$ for some $\psi: \R^d \to \R$, then we have $\varphi= -(-\psi)^*$ concave and upper-semicontinuous, since the Legendre transform is convex and lower-semicontinuous. Conversely, if $\varphi$ is concave and upper-semicontinuous, then we can define $\psi := -(-\varphi)^*$ and obtain
        $$
        \psi^c = -(-\psi)^* = -(-\varphi)^{**} = \varphi
        $$
        by Fenchel-Moreau theorem, since $-\varphi$ is convex and lower-semicontinuous.
        \item \textit{Metric Cost $c(x, y) = d(x, y)$} \\
        Let $X=Y$ with $c(x, y) = d(x, y)$, and let $\varphi: X \to \R$ be a measurable function. Then, 
        $$
        \varphi: X \to \R \ c\text{-concave} \iff \varphi \ 1\text{-Lipschitz, with } \varphi^c = -\varphi 
        $$
        Indeed, if $\varphi = \psi^c$ for some $\psi: X \to \R$ measurable, then we have
        \begin{align*}
            |\varphi(x_1) - \varphi(x_2)| 
            &= |\inf_{y\in X} \ (d(x_1, y) - \psi(y)) - \inf_{y\in X} \ (d(x_2, y) - \psi(y))| \\
            &\leq |\inf_{y\in X} \ (d(x_1, y) - \psi(y) - d(x_2, y) +\psi(y))| \\
            &\leq |\inf_{y\in X} \ d(x_1, x_2)| = d(x_1, x_2)
        \end{align*}
        by triangle inequality, and $\varphi$ is $1$-Lipschitz. \\
        Viceversa, suppose that $\varphi$ is $1$-Lipschitz, define $\psi:=-\varphi$ and obtain
        $$
        \psi^c(x)  = \inf_{y\in X} \ (d(x, y) + \varphi(y)) \leq d(x, x) + \varphi(x) = \varphi(x)
        $$
        and 
        $$
        \varphi(x) \leq \inf_{y\in X} \ (d(x, y) + \varphi(y)) = \inf_{y\in X} \ (d(x, y) - \psi(y)) = \psi^c(x) \quad \text{by Lipschitz}.
        $$
        In particular, the computation also shows that the $c$-transform of $\varphi$ is $\varphi^c = -\varphi$.
        \item \textit{Quadratic Cost $c(x, y) = \frac{1}{2}|x-y|^2$} \\
        Let $X = Y = \R^d$ with $c(x, y) = \frac{1}{2}|x-y|^2$, and let $\varphi: \R^d \to \R$ be a measurable function. Then,
        $$
        \varphi: \R^d \to \R \ c\text{-concave} \iff \frac{|x|^2}{2}-\varphi(x) \ \text{convex and LSC}.
        $$
        Note that the cost function $c(x, y) = \frac{1}{2} |x-y|^2$ has essentially the same structure of the cost $-x\cdot y$, since the interaction term between $x, y$ is the same. Indeed, the $c$-transform is given by
        \begin{align*}
            \varphi^c(y) 
            &= \inf_{x\in\R^d} \ \left(\frac{1}{2}|x-y|^2 - \varphi(x)\right) = \\
            &= \inf_{x\in\R^d} \ \left(\frac{1}{2}|x|^2 - x\cdot y + \frac{1}{2}|y|^2 - \varphi(x)\right) \\
            &= \frac{1}{2}|y|^2 + \inf_{x\in\R^d} \ \left(-x\cdot y + \frac{1}{2}|x|^2 - \varphi(x)\right) \\
            &= \frac{1}{2}|y^2| + \left(\frac{1}{2}|\cdot|^2 - \varphi\right)^*(y)
        \end{align*} 
        Then, one can prove similarly to $(i)$, that $\varphi$ is $c$-concave if and only if $\frac{1}{2}|\cdot|^2 - \varphi$ is convex and lower-semicontinuous, by Fenchel-Moreau theorem.
    \end{enumerate}
\end{example}

The previous example shows that we can interpret $c$-transforms as a generalization of the Legendre transform in $\R^d$ with quadratic cost $c(x, y) = -x\cdot y$. Now, we recall some results about the Legendre transform of a function and its biconjugate.

\begin{proposition}
    Let $\varphi: \R^d \to \R$ be a measurable function. Then, the Legendre transform $\varphi^*$ is convex and lower-semicontinuous.
\end{proposition}

\begin{theorem}
    (Fenchel-Moreau)
    Let $\varphi: \R^d \to \extRplus$ be a function. Then, the biconjugate $\varphi^{**}$ is the largest lower-semicontinuous convex function with $\varphi^{**} \leq \varphi$. \\ In particular, $\varphi^{**}=\varphi$ if and only if $\varphi$ convex and lower-semicontinuous.
\end{theorem}
\begin{proof}
    Theorem 12.2, \cite{rockafellar1997convex}.
\end{proof}

The following result generalizes the properties of the Legendre transform $\varphi^*$ and the biconjugate $\varphi^{**}$ to the case of $c$-transforms.

\begin{proposition}
    Let $c: X\times Y\to \R$ be a cost function, and $\varphi: X\to \R$ measurable. Then, it holds:
    \begin{enumerate}[(i)]
        \item $\varphi^c$ $c$-concave, upper-semicontinuous for $c$ continuous.
        \item $\varphi^{cc} \geq \varphi$, with equality if and only if $\varphi \ c\text{-concave}$.
        Moreover, $\varphi^{cc}$ is smallest $c$-concave function larger than $\varphi$, called $c$-concavification of $\varphi$.
    \end{enumerate}
\end{proposition}
\begin{proof} \ \\
    \textit{(i)} Note that if the cost function $c$ is continuous, then the map $y \mapsto c(x, y) - \varphi(x)$ is continuous for all $x\in X$. Therefore, the map $$\varphi^c(y) = \inf_{x\in X} \ (c(x, y) - \varphi(y))$$ is upper-semicontinuous, as a infimum of continuous functions.\\
    \textit{(ii)} By definition of $c$-transform we have
    \begin{align*}
        \varphi^{cc}(x)
        &= \inf_{y\in Y} \ (c(x, y) - \varphi^c(y)) \\
        &= \inf_{y\in Y} \ (c(x, y) - \inf_{\tilde x\in X} \ (c(\tilde x, y) - \varphi(\tilde x))) \\
        &= \inf_{y\in Y} \ \underbrace{\sup_{\tilde x\in X} (c(x, y) - c(\tilde x, y) + \varphi(\tilde x))}_{\geq c(x, y) - c(x, y) + \varphi(x) = \varphi(x)} \\
        &\geq \inf_{y\in Y} \ \varphi(x) = \varphi(x)
    \end{align*}
    Now, note that $\varphi^{ccc} = \varphi$, since
    \begin{align*}
        \varphi^{ccc}(y) 
        &= \inf_{x\in X} \ (c(x, y) - \varphi^{cc}(x)) \\
        &= \inf_{x\in X} \ (c(x, y) - \inf_{\tilde y\in Y} \ \sup_{\tilde x\in X} (c(x, \tilde y) - c(\tilde x, \tilde y) + \varphi(\tilde x))) \\
        &= \inf_{x\in X} \ \sup_{\tilde y\in Y} \ \inf_{\tilde x \in X} \ (c(x, y) - c(x, \tilde y) + c(\tilde x, \tilde y) - \varphi(\tilde x))
    \end{align*}
    Then, by choosing $\tilde y = y$ we get
    \begin{align*}
        \varphi^{ccc}(y) 
        &=\inf_{x\in X} \ \sup_{\tilde y\in Y} \ \inf_{\tilde x \in X} \ (c(x, y) - c(x, \tilde y) + c(\tilde x, \tilde y) - \varphi(\tilde x)) \\
        &\geq \inf_{x\in X} \ \inf_{\tilde x \in X} \ (c(x, y) - c(x, y) + c(\tilde x, y) - \varphi(\tilde x)) \\
        &= \inf_{\tilde x\in X} \ (c(\tilde x, y) - \varphi(\tilde x)) = \varphi^c(y)
    \end{align*}
    Conversely, by choosing $\tilde x = x$ we obtain $\varphi^{ccc} \leq \varphi$, and conclude that $\varphi^{ccc} = \varphi$. \\
    Note that $\varphi^{cc}$ is $c$-concave by definition. Viceversa, if $\varphi$ $c$-concave, then $\varphi = \psi^{c}$ for some $\psi: Y \to \R$, and we have $\varphi^{cc} = (\psi^c)^{cc} = \psi^c = \varphi$ by previous computation.
\end{proof}

Finally, we state the duality theorem in its most general form.

\begin{theorem}
    (Duality, II) Let $c: X\times Y\to \extRplus$ be a lower-semicontinuous and bounded below cost function for a Kantorovich problem with marginals $\mu\in\mathcal{P}(X), \ \nu\in\mathcal{P}(Y)$. Then, 
    \begin{align*}
        \inf_{\pi\in Cpl(\mu, \nu)} \ \int_{X\times Y} c \ d\pi 
        &= \sup_{\substack{(\varphi, \psi)\in C_b(X)\times C_b(Y) \\ c\geq \varphi\oplus\psi}} \left(\int_X \varphi \ d\mu \ + \ \int_Y \psi \ d\nu \right) \\
        &= \sup_{\substack{(\varphi, \psi)\in L^1(\mu)\times L^1(\nu) \\ c\geq \varphi\oplus\psi}} \left(\int_X \varphi \ d\mu \ + \ \int_Y \psi \ d\nu \right) \\
        &= \sup_{\psi\in L^1(\nu)} \left(\int_X \psi^c \ d\mu + \int_Y \psi \ d\nu\right) \\
        &= \sup_{\varphi\in L^1(\mu)} \left(\int_X \varphi \ d\mu + \int_Y \varphi^c \ d\nu\right)
    \end{align*}
    Moreover, if $\varphi^*\in L^1(\mu), \psi^*\in L^1(\nu)$ be optimizer of the dual problem, then
    $$
    \pi\in Cpl(\mu, \nu) \text{ optimal} \iff c = \varphi^* \oplus \psi^* \quad \pi-a.s.
    $$
    Similarly, if $\pi^*$ is an optimizer of the Kantorovich problem, then
    $$
    \varphi\in L^1(\mu), \ \psi\in L^1(\nu) \text{ s.t. $\varphi\oplus\psi\leq c$ optimal } \iff c = \varphi \oplus \psi \quad \pi^*-a.s.
    $$
\end{theorem}
\begin{proof}
    The thesis follows by Theorem \ref{thm:strong_duality}, by noting that for any $\varphi: X \to \R, \ \psi: Y\to\R$ it holds $\varphi + \varphi^c\leq c$ by definition of $c$-transform, with 
    $$
    \int_X \varphi \ d\mu + \int_Y \psi \ d\nu \leq \int_X \varphi \ d\mu + \int_Y \varphi^c \ d\nu.
    $$
\end{proof}

A particular case of the duality theorem is the Kantorovich-Rubenstein formula.

\begin{corollary}
    \label{thm:Kantorovich-Rubenstein_classicOT}
    (Kantorovich-Rubenstein) 
    The Kantorovich problem with metric cost $c(x, y) = d(x, y)$ on a Polish space $X$, and marginals $\mu, \nu\in\mathcal{P}_1(X)$ satisfies the duality
    $$
    \inf_{\pi\in Cpl(\mu, \nu)} \int_{X\times Y} c \ d\pi = \sup_{\varphi  \ 1-\text{Lip}} \left(\int_X \varphi \ d \mu - \int_X \varphi \ d \nu\right).
    $$
    Moreover, the problem both problems admit optimizers. In particular, if $\varphi^*$ $1$-Lipschitz optimizer, then 
    $$
    \pi\in Cpl(\mu, \nu) \text{ optimal} \iff \varphi^*(x)-\varphi^*(y)=d(x, y) \text{ for } \pi-a.s. \ x, y\in X.
    $$
\end{corollary}
\begin{proof}
    Note that the cost function $c(x, y) = d(x, y)$ is continuous, then we get by duality theorem 
    \begin{align*}
        \inf_{\pi\in Cpl(\mu, \nu)} \int c \ d\pi 
        &= \sup_{\varphi\in L^1(\mu)} \left(\int_X \varphi \ d\nu + \int_Y \varphi^c \ d\nu\right) \\
        &= \sup_{\varphi  \ 1-\text{Lip}} \left(\int_X \varphi \ d\nu - \int_Y \varphi \ d\nu\right)
    \end{align*}    
    since the transform of a measurable map $\varphi$ is $\varphi^c=-\varphi$, and the $c$-concave functions are exactly the Lipschitz functions of constant $1$.
\end{proof}

\subsection{Structural Results}
Now, we want to better study the structure of the space of optimizers of the Kantorovich problem.

\begin{definition}
    A set $\Gamma\subseteq X\times Y$ is cyclically monotone respect to a cost function $c: X\times Y\to \extRplus$, or also $c$-cyclically monotone ($c-cm$), if 
    $$
    \sum_{i=1}^n c(x_i, y_i) \leq \sum_{i=1}^n c(x_i, y_{i+1})
    $$
    for all finite sequences $(x_1, y_1), \dots, (x_n, y_n)\in\Gamma$. \\
    Equivalently, $\Gamma\subseteq X \times Y$ is $c$-cyclically monotone if 
    $$
    \sum_{i=1}^n c(x_i, y_i) \leq \sum_{i=1}^n c(x_i, y_{P(i)})
    $$
    for all sequences $(x_1, y_1), \dots, (x_n, y_n)\in\Gamma$, and all indices permutations $P\in Sym(n)$.
\end{definition}

\begin{remark}
    The cyclically monotonicity is a geometric condition on the support of couplings, and is independent of the distribution of mass given by the coupling.
\end{remark}


\begin{definition}
    A probability measure $\pi\in\mathcal{P}(X\times Y)$ is $c$-cyclically monotone for a cost $c:X\times Y\to\extRplus$ if $\pi(\Gamma)=1$ for some $c$-cyclically monotone set $\Gamma\subseteq X\times Y$.
\end{definition}

\begin{remark}
    In particular, if the support $supp(\pi)\subseteq X\times Y$ of $\pi$ is $c$-cyclically monotone, then $\pi$ is $c$-cyclically monotone. We will see in the following, that the geometric structure of $supp(\pi)$ is determinant for the optimality of the coupling.
\end{remark}

\begin{proposition}
    \label{lemma:optimal_implies_cm}
    Let $c: X\times Y \to [0, \infty]$ be a non-negative and continuous and bounded below cost function for a Kantorovich problem with marginals $\mu\in\mathcal{P}(X), \nu\in\mathcal{P}(Y)$. If $\pi^* \in Cpl(\mu, \nu)$ optimal coupling with finite value, then $supp(\pi)\subseteq X\times Y$ $c$-cyclically monotone. 
\end{proposition}
\begin{proof}
    Assume by contradiction that $$\sum_{i=1}^n c(x_i, y_i) > \sum_{i=1}^n c(x_i, y_{i+1})$$ for some $(x_1, y_1), \dots, (x_n, y_n)\in supp(\pi)$. Then, there exists $U_1, \dots, U_n$ open disjoint neighborhoods of $x_1, \dots, x_n$, and $V_1, \dots, V_n$ open disjoint neighborhoods of $y_1, \dots, y_n$ such that 
    \begin{equation}
    \label{eq:contraddiction_c_cyclical_monotone}
        \sum_{i=1}^n c(u_i, v_i) > \sum_{i=1}^n c(u_i, v_{i+1})
    \end{equation}
    for all $u_i\in U_i, v_i\in V_i$, by continuity of $c$. 
    Now, consider the probability measures $$\pi_i := \frac{1}{m_i} \pi^*_{|U_i\times V_i} \quad \text{for $i=1, \dots, n$},$$ where $\pi_{|U_i\times V_i}(A) = \pi^*(A \cap (U_i \times V_i))$ and $m_i = \pi^*(U_i \times V_i)$. \\
    Finally, let $\mu_i, \nu_i$ be the marginals of $\pi_i$, and consider the probability measure
    $$
    \tilde\pi := \pi^* + \frac{\min_{i=1, \dots, n} m_i}{n} \left(\sum_{i=1}^n \mu_i\otimes \nu_{i+1} \ - \ \pi_i \right).
    $$
    Note that 
    $$
    \pi^* - \frac{\min_{i=1,\dots,n} m_i}{n} \sum_{i=1}^n \pi_i \in Cpl\left(\mu - \frac{\min_{i=1,\dots,n} m_i}{n}\sum_{i=1}^n \mu_i, \ \nu - \frac{\min_{i=1,\dots,n} m_i}{n}\sum_{i=1}^n \nu_i\right)
    $$
    by Lemma \ref{lemma:coupling_composition}. Then, we have
    $$
    \tilde \pi = \pi^* - \frac{\min_{i=1,\dots,n} m_i}{n}\sum_{i=1}^n \pi_i \ + \ \frac{\min_{i=1,\dots,n} m_i}{n} \sum_{i=1}^n \mu_i\otimes\nu_{i+1} \in Cpl(\mu, \nu). 
    $$
    Finally, one can prove that the coupling $\tilde \pi\in Cpl(\mu, \nu)$ achieves a lower average cost than the optimal coupling $\pi\in Cpl(\mu, \nu)$ using $\ref{eq:contraddiction_c_cyclical_monotone}$.
\end{proof}

Now, we introduce another geometric notion which will be related to the structure of the support of optimal couplings.

\begin{definition}
    The $c$-superdifferential of a $c$-concave function $\varphi : X \to \R$ is the set $$\partial^c \varphi := \{(x, y)\in  X \times Y \mid \varphi(x) + \varphi^c(y) = c(x, y)\}.$$
    In particular, the $c$-superdifferential of $\varphi$ at $x\in X$ is the set $$\partial^c\varphi(x) = \{y \in Y \mid (x, y)\in \partial^c\varphi\}.$$
\end{definition}

\begin{remark}
    Intuitively, the $c$-superdifferential of a $c$-concave function $\varphi$ contains the points satisfying the optimality of the dual problem.\\Indeed, we have see in the duality theorem that, if $\varphi$ optimizer of the dual problem, then a coupling $\pi\in Cpl(\mu, \nu)$ is optimal if and only if 
    $$
    1= \pi(\{(x, y)\in X\times Y | \varphi(x) + \varphi^c(y) = c(x, y)\}) = \pi(\partial^c\varphi).
    $$
    In particular, if any coupling $\pi$ concentrated on the $c$-superdifferential $\partial^c \varphi$ of a $c$-concave function $\varphi$ is optimal for the Kantorovich problem.
\end{remark}

\begin{theorem}
    \label{thm:Rockafellar}
    (Rockafellar) 
    Let $c: X\times Y \to \R$ be a cost function, then 
    $$\Gamma\subseteq X\times Y \ c - \text{cyclically monotone} \iff \Gamma \subseteq \partial^c \varphi \text{ for some $c$-concave function $\varphi$}$$
\end{theorem}
\begin{proof}
    Theorem 24.8, \cite{rockafellar1997convex}.
\end{proof}

\begin{theorem}
    \label{thm:fundamental_OT}
    (Fundamental OT Theorem) 
    Let $c: X \times Y \to \R$ be a lower-semicontinuous cost function such that $|c(x, y)| \leq a(x) + b(y)$ for some $a\in L^1(\mu), b\in L^1(\nu)$, and let $\pi\in Cpl(\mu, \nu)$ be a coupling of $\mu\in \mathcal{P}(X), \nu\in\mathcal{P}(Y)$ with finite value. \\
    Then, the following are equivalent:
    \begin{enumerate}[(i)]
        \item $\pi$ optimizer of the Kantorovich problem $P^K_c(\mu, \nu)$.
        \item $supp(\pi)\subseteq X\times Y$ $c$-cyclically monotone.
        \item $supp(\pi)\subseteq \partial^c \varphi$ for some $c$-concave function $\varphi$.
    \end{enumerate}
    Moreover, duality holds and $\varphi\in L^1(\mu), \varphi^c\in L^1(\nu)$ are the optimizers of the dual problem.
\end{theorem}
\begin{proof}
    Firstly, we assume that the cost $c: X\times Y \to \R$ is a non-negative continuous function; otherwise we can replace $c(x, y)$ by $\tilde c(x, y) = c(x, y) + a(x) + b(y)$. \\
    Then, note that the implications $\textit{(i)}\implies\textit{(ii)}$ and $\textit{(ii)}\implies\textit{(iii)}$ are guaranteed by Lemma \ref{lemma:optimal_implies_cm}, and Theorem \ref{thm:Rockafellar} respectively. \\
    Now, suppose that $supp(\pi)\subseteq \partial^c \varphi$ for some $c$-concave function $\varphi$, and we prove that the measure $\pi$ is an optimal coupling. Let $\psi := \varphi^c$ be the $c$-transform of $\varphi$, and note that by definition of $c$-transform we have
    $$
    |\varphi\oplus\psi| 
    = |\varphi\oplus\varphi^c| 
    \leq |c| \leq a\oplus b.
    $$
    Then, $\varphi_+ \in L^1(\mu), \ \psi_+\in L^1(\nu)$ since $a\in L^1(\mu), \ b\in L^1(\nu)$, and therefore the integrals $\int_X \varphi \ d\mu, \int_Y \psi \ d\nu$ are well-defined. \\
    Now, consider the increasing sequences $\varphi_n := (\varphi \vee (-n))\wedge n, \ \psi_n := (\psi \vee (-n))\wedge n$, and note that
    $$
    \varphi_n \to \varphi, \ \psi_n\to \psi \quad \text{for } n\to\infty.
    $$
    Moreover, $\varphi_n\oplus\psi_n\leq c$ and hence $\varphi_n \in L^1(\mu), \ \psi_n\in L^1(\nu)$, since $\varphi_n, \psi_n$ bounded in $[-n,n]$. Then, we have
    \begin{align*}
        \int_{X\times Y} c \ d\pi 
        &= \int_{X\times Y} \varphi\oplus\psi \ d\pi \quad \text{since } supp(\pi)\subseteq \partial^c\varphi \\
        &= \lim_{n\to\infty}\int_{X\times Y} \varphi_n\oplus\psi_n \ d\pi \quad \text{by monotone convergence} \\
        &= \lim_{n\to\infty}\int_X \varphi_n \ d\mu + \int_Y \psi_n \ d\nu \quad \text{since } \varphi_n\in L^1(\mu), \psi_n\in L^1(\nu) \ \\
        &= \int_X \varphi \ d\mu + \int_Y \psi \ d\nu. 
    \end{align*}
    Thus, $\varphi\in L^1(\mu), \ \psi\in L^1(\nu)$ are optimizers of the dual problem, and the coupling $\pi$ is an optimizer of the Kantorovich problem. \\
    More generally, we can approximate the cost function $c$ with a non-decreasing sequence of continuous functions and prove the result similarly.
\end{proof}

\begin{remark}
    The fundamental theorem of optimal transport shows that the optimality of a coupling depends only on the geometry of the support relatively to the cost function, and does not depend on the distribution of mass on the support. For instance, if $\pi$ is an optimal coupling, and $\tilde\pi$ is another coupling with $supp(\tilde\pi)\subseteq supp(\pi)$, then $\tilde\pi$ is also optimal, because the $c$-cyclical monotonicity is also inherited by subsets.
\end{remark}

\begin{remark}
    A particular case is realized when the $c$-superdifferential is $$\partial^c\varphi(x) =\{*\} \quad \text{for } \mu-a.e. \ x\in X.$$
    In this case, the optimal couplings $\pi\in Cpl(\mu, \nu)$ have their support $supp(\pi)$ contained in a space $\partial^c\varphi$ with only one degree of freedom. In other words, optimal couplings are supported on the graph of a function, i.e. the coupling is induced by a transport map.
\end{remark}

Now, we will analyze some results about special cases of the classical optimal transport problem, namely the real case $X=Y=\R^d$ with quadratic cost.

\begin{definition}
    A function $T: \R^d \to \R^d$ is monotone respect to a probability measure $\mu\in\mathcal{P}(\R^d)$ if $T = \nabla\varphi$ $\mu$-a.s. for some convex function $\varphi: \R^d \to \R$.
\end{definition}

\begin{theorem}
    (Rademacher) 
    Let $f: \R^d \to \R$ be a convex function. Then, $f$ is a.e. differentiable.
\end{theorem}

\begin{theorem}
    The Kantorovich problem with quadratic cost $c(x, y) = x\cdot y$ on $\R^d$, and marginals $\mu, \nu\in\mathcal{P}_2(\R^d)$ s.t. $\mu\ll\lambda$ admits a unique optimal coupling $\pi^*\in Cpl(\mu, \nu)$ of the form $\pi^* = (Id, \nabla\varphi)\mu$ for some convex function $\varphi: \R^d \to \R$. In particular, there exists a unique optimal transport map $T: \R^n \to \R^n$ from $\mu$ to $\nu$ of the form $T=\nabla\varphi$ for some convex function $\varphi: \R^n\to\R$.    
\end{theorem}
\begin{proof}
    Note that the existence of an optimizer is guaranteed by the continuity of the quadratic cost $c(x, y) = x\cdot y$, by Theorem \ref{thm:optimal_plan_existence}. Moreover, $c$ is bounded as
    $$
    |c(x,y)| = |x\cdot y| \leq \frac{1}{2}|x|^2 + \frac{1}{2}|y|^2
    $$
    with $a(x)=\frac{1}{2}|x|^2\in L^1(\mu), \ b(y)=\frac{1}{2}|y|^2\in L^1(\nu)$, since $\mu$ absolutely continuous respect to the Lebesgue measure. Then, a coupling $\pi\in Cpl(\mu, \nu)$ is optimal if and only if $supp(\pi)\subseteq \partial^c \varphi$ for some $\varphi:\R^d \to \R$ $c$-concave function, by Theorem \ref{thm:fundamental_OT}. \\
    Recall that, for the quadratic cost $c(x, y)=x\cdot y$, a measurable map $\varphi: \R^d\to\R$ we have
    $$
    \varphi \ c\text{-concave} \iff \varphi \text{ convex and LSC}
    $$
    Therefore, a coupling $\pi\in Cpl(\mu, \nu)$ is optimal if and only if $supp(\pi)\subseteq \partial^c\varphi$ for some $\varphi:\R^d\to \R$ convex and lower-semicontinuous. In particular, $\varphi$ is $\mu$-a.e. differentiable by Rademacher theorem, and the gradient $\nabla\varphi$ is $\mu$-a.e. well-defined on $\R^d$, and consequently
    $$
    \nabla\varphi(x) = y \iff y\in\partial^c\varphi(x) \quad \text{for } \mu-a.e. \ x\in\R^d.
    $$
    Thus, the $c$-superdifferential is a singleton $\mu$-a.e, and any optimal coupling $\pi$ is supported on the graph of the gradient of differentiable functions. \\
    Now, suppose that $\pi_1, \pi_2$ are optimal couplings supported on the graph of $T_1, T_2$ respectively. Then, the convex combination $\pi_3 = \frac{1}{2}(\pi_1 + \pi_2)$ is also optimal by linearity, and thus is supported on the graph of a function. Consequently, $T_1 = T_2$ and $\pi_1 = \pi_2$. \\
    At last, note that any transport map between $\mu$ and $\nu$ of the type $T=\nabla\varphi$ is optimal, since its graph is $c$-cyclically monotone.
\end{proof}

\begin{theorem}
    (Brenier)
    The Kantorovich problem with quadratic cost $c(x, y) = \frac{1}{2}|x-y|^2$ on $\R^d$, and marginals $\mu, \nu\in\mathcal{P}_2(\R^d)$ s.t. $\mu\ll\lambda$ admits a unique optimal coupling $\pi^*\in Cpl(\mu, \nu)$ of the form $\pi^* = (Id, \nabla\varphi)\mu$ for some convex function $\varphi: \R^d \to \R$. In particular, there exists a unique optimal transport map $T: \R^n \to \R^n$ from $\mu$ to $\nu$ of the form $T=\nabla\varphi$ for some convex function $\varphi: \R^n\to\R$.    
\end{theorem}
\begin{proof}
    The thesis follows by observing that 
    \begin{align*}
        \inf_{\pi\in Cpl(\mu, \nu)} \int \frac{1}{2}|x-y|^2 \ d\pi
        &= \inf_{\pi\in Cpl(\mu, \nu)} \int \frac{1}{2} |x|^2 - x\cdot y + \frac{1}{2}|y|^2  \ d\pi \\
        &= \inf_{\pi\in Cpl(\mu, \nu)} \left(\int |x|^2 \ d\mu + \int |y|^2 \ d\nu\right) + \int -x\cdot y \ d\pi \\
        &= \frac{1}{2}\left(\int |x|^2 \ d\mu + \int |y|^2 \ d\nu\right) + \inf_{\pi\in Cpl(\mu, \nu)}  \int -x\cdot y \ d\pi.
    \end{align*}
    Then, the minimization problem for the quadratic cost $c(x, y) = \frac{1}{2}|x-y|^2$ is reduced to the minimization problem for the quadratic cost $c(x, y) = -x\cdot y$.
\end{proof}

\subsection{Wasserstein Space}
In this section, we introduce a metric on a subspace of the space $\mathcal{P}(X)$ of probability measures on a Polish space $X$, namely the space of $p$-finite moment probability measures.

\begin{definition}
    A probability measure $\mu$ of a Polish space $X$ has $p$-finite moment, with $p\geq 1$, if $$\int_X d(x, x_0)^p \ d\mu(x)<\infty$$ for some $x_0\in X$.
\end{definition}
We will denote the space of finite $p$-moment probability measures on $X$ by $\mathcal{P}_p(X)$.

\begin{remark}
    The notion is well defined, in the sense that is does not depend on the choice of the point $x_0\in X$. Indeed, for all $x_0, x_1\in X$ we have
    $$
    d(x, x_1)^p 
    \leq (d(x, x_0) + d(x_0, x_1))^p 
    \leq 2^{p-1} (d(x, x_0)^p + d(x_0, x_1)^p)
    $$
    since $(a+b)^p\leq 2^{p-1}(a^p+b^p)$ for all $a,b\geq 0, \ p\geq 1$ by Jensen inequality with $t\mapsto t^p$. \\
    Therefore, if $$\int_X d(x, x_0)^p \ d\mu(x)<\infty$$ for some $x_0\in X$, and if $x_1\in X$ is a distinct point, we get 
    $$
    \int_X d(x, x_1)^p \ d\mu(x)
    \leq 2^{p-1} \left(\int_X (d(x, x_0)^p \ d\mu(x) + d(x_0, x_1)^p\right) < \infty 
    $$
\end{remark}

\begin{definition}
    The $p$-Wasserstein distance, with $p\geq 1$, between $\mu\in\mathcal{P}_p(X)$ and $\nu\in\mathcal{P}_p(Y)$ is given by
    $$
    \mathcal{W}_p(\mu, \nu) 
    := \left(\inf_{\pi\in Cpl(\mu, \nu)} \int_{X\times Y} d(x, y)^p \ d\pi(x, y) \right)^{1/p}
    = \left(\inf_{X\sim\mu, Y\sim\nu} \E[d(X, Y)^p]\right)^{1/p}
    $$
\end{definition}

\begin{remark}
    The $p$-Wasserstein distance between $\mu$ and $\nu$ is a Kantorovich optimal transport problem with $X=Y$ and metric cost $d(x, y)^p$. Hence, we have $$\mathcal{W}_p(\mu, \nu) = P^K_{d^p}(\mu, \nu)^{1/p}. $$
    In particular, the optimization problem defined by $\mathcal{W}_p(\mu, \nu)$ admits minimizers by continuity of the metric cost $c(x, y) = d(x, y)$.
\end{remark}

\noindent
We will prove in the following that $\mathcal{W}_p$ is well defined on $\mathcal{P}_p(X)$ in the sense that is assumes only finite values, and moreover we will prove that $\mathcal{W}_p$ is a metric, for any $p\geq 1$.

\begin{lemma}
    (Gluing)
    \label{lemma:gluing_lemma} Let $\mu_1\in \mathcal{P}(X_1), \ \mu_2\in \mathcal{P}(X_2), \ \mu_3\in \mathcal{P}(X_3)$ be probabilities on the Polish spaces $X_1, X_2, X_3$, and consider the couplings $\pi_{12}\in Cpl(\mu_1, \mu_2), \ \pi_{23}\in Cpl(\mu_2, \mu_3)$. \\
    Then, there exists a probability distribution $\pi\in\mathcal{P}(X_1 \times X_2 \times X_3)$ such that $p_{12}\pi = \pi_{12}$ and $p_{23}\pi = \pi_{23}$, where $p_{ij}: X_1\times X_2 \times X_3 \to X_i\times X_j$ projection.
\end{lemma}
\begin{proof}
    Note that the couplings $\pi_{12}\in Cpl(\mu_1, \mu_2), \ \pi_{23}\in Cpl(\mu_2, \mu_3)$ can be disintegrated along their common marginal $\mu_2$ as
    \begin{align*}
        \pi_{12}(dx_1, dx_2) &= \pi_{12}(dx_1 \mid x_2) \ \mu_2(dx_2) \\
        \pi_{12}(dx_2, dx_3) &= \pi_{23}(dx_3 \mid x_2) \ \mu_2(dx_2)
    \end{align*}
    Then, we can construct the $3$-marginal probability measure $\pi$ as 
    $$
    \pi(dx_1, dx_2, dx_3) := \pi_{12}(dx_1 \mid x_2) \ \mu_2(dx_2) \ \pi_{23}(dx_3 \mid x_2).
    $$
    Hence, $\pi\in \mathcal{P}(X_1 \times X_2 \times X_3)$ has projections $\pi_{12}, \pi_{23}$ by construction.
\end{proof}

\begin{theorem}
    $(\mathcal{P}_p(X), \ \mathcal{W}_p)$ metric space, called $p$-Wasserstein space, for $p\geq 1$.
\end{theorem}
\begin{proof} \ \\
    \textit{Well-def.)} Let $\mu, \nu\in \mathcal{P}_p(X)$, let $\pi\in Cpl(\mu, \nu)$, and fix a point $x_0\in X$. Thus, we get
    \begin{align*}
        \mathcal{W}_p(\mu, \nu)^p 
        &\leq \int d(x, y)^p \ d\pi(x, y) \\
        &\leq 2^{p-1}\int d(x, x_0)^p \ d\pi(x, y) + 2^{p-1}\int d(x_0, y)^p \ d\pi(x, y) \\
        &= 2^{p-1}\int d(x, x_0)^p \ d\mu(x) + 2^{p-1}\int d(x_0, y)^p \ d\nu(y) < +\infty,
    \end{align*}
    since $\mu, \nu$ have finite $p$-moment. \\
    \textit{Symmetry)} Let $\pi\in Cpl(\mu, \nu)$, and consider the transposed coupling $\pi^T := T\pi$ obtained as pushforward of $\pi$ under the map $T(x, y) = T(y, x)$. Note, that $\pi^T\in Cpl(\nu, \mu)$ since
    $$
    p_1\pi^T = p_1T\pi = p_2\pi = \nu, \quad p_2\pi^T = p_2 T\pi = p_1\pi = \mu.
    $$
    Thus, we have
    \begin{align*}
        \int d(x, y)^p  \ d \pi(x, y) 
        &= \int d(T(x, y))^p \ d \pi(x, y) \quad \text{by symmetry of $d$} \\
        &= \int d(x, y)^p \ d\pi^T(x, y) \quad \text{by change of variable, since } T\pi = \pi^T.
    \end{align*}
    Therefore every coupling of $(\mu, \nu)$ gives a coupling of $(\nu, \mu)$ with exactly the same cost, and by taking infimum over all couplings, we get
    $$
    \inf_{\pi\in Cpl(\mu, \nu)} \int_{X\times Y} d(x, y)^p \ d\pi(x, y)
    \leq
    \inf_{\pi\in Cpl(\nu, \mu)} \int_{X\times Y} d(x, y)^p \ d\pi(x, y).
    $$
    The symmetry follows by changing the roles of $\mu, \nu$ in the inequality. \\
    \textit{Non-Negativity)} Let $\mu\in\mathcal{P}_p(X)$ and note that the diagonal coupling $\pi:=(Id, Id)\pi$ has cost
    $$
    \int d(x, y)^p \ d\pi(x, y) = \int d(x, x)^p \ d\mu(x) = 0 \quad \text{by change of variable}
    $$
    since $d=0$ on the diagonal of $X\times X$. Then, we have
    $$
    0 \leq \mathcal{W}_p(\mu, \mu)^p = \inf_{\pi\in Cpl(\mu, \mu)} \int d(x, y)^p \ d\pi(x, y) \leq 0.
    $$
    Conversely, suppose that 
    $$
    \mathcal{W}_p(\mu, \nu)^p = \inf_{\pi\in Cpl(\mu, \nu)} \int d(x, y)^p \ d\pi(x, y) = 0
    $$
    for some $\mu, \nu\in\mathcal{P}_p(X)$. Hence, $d(x, y) = 0$ $\pi$-a.s. for some $\pi\in Cpl(\mu, \nu)$, which implies that $x=y$ $\pi$-a.e., since $d$ non-negative. Therefore, $\pi$ is entirely concentrated on the diagonal of $X\times X$, i.e. $\pi$ induced by identity map and $\nu = (Id, Id)\mu = \mu$. \\
    \textit{Triangle inequality)} Let $\mu_1, \mu_2, \mu_3\in\mathcal{P}_p(X)$, and let $\pi_{12}\in Cpl(\mu_1, \mu_2), \ \pi_{23}\in Cpl(\mu_2, \mu_3)$ be optimal coupling for the respective Wasserstein problems. \\
    Then, there exists a probability distribution $\pi\in\mathcal{P}(X_1 \times X_2 \times X_3)$ such that $p_{12}\pi = \pi_{12}$ and $p_{23}\pi = \pi_{23}$, by gluing lemma. Now, let $(Z_1, Z_2, Z_3)$ be a random vector with law $\pi$, such that $(Z_1, Z_2)\sim\pi_{12}$ and $(Z_2, Z_3)\sim\pi_{23}$ by construction. Finally, we get
    \begin{align*}
        \mathcal{W}_p(\mu_1, \mu_3) 
        &= \left(\inf_{X\sim\mu_1, Y\sim\nu_3} \E[d(X, Y)^p]\right)^{1/p} \\
        &\leq \E[d(Z_1, Z_3)^p]^{1/p} \quad \text{since } Z_1\sim\mu_1, \ Z_3\sim\mu_3\\
        &\leq  \E[d(Z_1, Z_2)^p]^{1/p} + \E[d(Z_2, Z_3)^p]^{1/p} \quad \text{by Minkowski inequality} \\
        &= \mathcal{W}_p(\mu_1, \mu_2) + \mathcal{W}_p(\mu_2, \mu_3),
    \end{align*} 
    since $\pi_{12}\sim (Z_1, Z_2), \ \pi_{23}\sim (Z_2, Z_3)$ optimal couplings.
\end{proof}

\begin{remark}
    For $p=1$, we have by Kantorovich-Rubenstein formula 
    $$
    \mathcal{W}_1(\mu, \nu) = \inf_{\pi\in Cpl(\mu, \nu)} \int d(x, y) \ d\pi(x, y) = \sup_{\varphi \ 1\text{-Lip}} \int \varphi \ d(\mu - \nu).
    $$
    The Wasserstein distance $\mathcal{W}_1$ of order $1$ is sometimes called the Kantorovich-Rubenstein distance on $\mathcal{P}_1(X)$.
\end{remark}

\begin{example}
    The $p$-Wasserstein distance between Dirac measures coincides with the usual metric for any $p\geq 1$. Indeed, the space of couplings $Cpl(\delta_x, \delta_y)$ consists only of the independent coupling $\delta_{(x, y)} = \delta_x \otimes \delta_y$ by Example \ref{example:dirac_coupling}, and we get
    $$
    \mathcal{W}_p(\delta_{x}, \delta_y) = \left(\int d(x, y) \ d(\delta_{x}\otimes\delta_y)(x, y)\right)^{1/p} = d(x, y).
    $$
\end{example}

\begin{proposition}
    \label{prop:wasserstein_space_ordering}
    $\mathcal{W}_p\leq \mathcal{W}_q$ for $1\leq p \leq q<+\infty$.
\end{proposition}
\begin{proof}
    Let $r:= q/p$ and $r':=q/(q-p)$ such that $r, r'\geq 1$, $1/r + 1/r' = 1$. Then, by Holder inequality we get
    \begin{align*}
        ||d(X, Y)^p \cdot 1||_1 
        \leq ||d(X, Y)^p||_r \cdot ||1||_{r'}
        = \E[d(X, Y)^{rp}]^{1/r}
        = \E[d(X, Y)^q]^{p/q}.
    \end{align*}
    Then, we have
    \begin{align*}
        \mathcal{W}_p(\mu, \nu)^p = \inf_{X\sim\mu, \ Y\sim\nu} \E[d(X, Y)^p] \leq \inf_{X\sim\mu, Y\sim\nu} \E[d(X, Y)^q]^{p/q} = \mathcal{W}_q(\mu, \nu)^p. 
    \end{align*}
\end{proof}

\begin{remark}
    Up to now, we have introduced two a priori distinct topologies on the space $\mathcal{P}_p(X)$:
    \begin{itemize}
        \item The weak topology induced by weak convergence on the subspace $\mathcal{P}_p(X)\subseteq \mathcal{P}(X)$
        \item The metric topology induced by the $p$-Wasserstein distance
    \end{itemize}
We will now investigate the relation between these two notions of convergence.
\end{remark}

\begin{theorem}
    (Skorohod representation) Let $\{\mu_n\}\subseteq\mathcal{P}(X)$ be a sequence of probability measures on a Polish space $X$ converging weakly to a probability $\mu\in\mathcal{P}(X)$. \\
    Then, there exist a sequence $\{Z_n\}_n$ of $X$-valued random variables on a probability space $(\Omega, \mathcal{F},\Pbb)$ such that
    \begin{enumerate}[(i)]
        \item $X_n \to X$ almost surely for $n\to\infty$.
        \item $X_n\sim\mu_n, \ X\sim\mu$ for all $n\in\N$.
    \end{enumerate}
\end{theorem}

\begin{theorem}
    Let $X$ be a bounded Polish space. Then, $\mathcal{W}_p$ metrizes the weak topology of $\mathcal{P}(X)$, for $p\geq 1$, i.e. $$\mu_n\to\mu \text{ weakly in } \mathcal{P}(X) \iff \mathcal{W}_p(\mu_n, \mu) \to 0.$$
\end{theorem}
\begin{proof}
    Let $\{\mu_n\}\subseteq\mathcal{P}(X)$ be a sequence converging weakly to $\mu\in\mathcal{P}(X)$. Then, there exist $\{Z_n\}_n, Z$ random variables with laws $\mu_n, \mu$ respectively, such that $$Z_n\to Z \quad \text{$\Pbb$-almost surely for $n\to\infty$.}$$ by Skorohod representation theorem. Let $\pi_n \in Cpl(\mu_n, \mu)$ be the law of the vector $(Z_n, Z)$, and note that
    \begin{align*}
        \limsup_{n\to\infty} \mathcal{W}_p(\mu_n, \mu)^p 
        &= \limsup_{n\to\infty} \ \inf_{\pi\in Cpl(\mu_n, \nu)} \int d(x, y) \ d\pi(x, y) \\
        &\leq \limsup_{n\to\infty} \int d(x, y) \ d\pi_n(x, y) \\
        &= \limsup_{n\to \infty} \E[d(Z_n, Z)] \quad \text{since } \pi_n\sim(Z_n, Z) \\
        &=0 \quad \text{by DCT.}
    \end{align*}
    Conversely, suppose that $\mathcal{W}_p(\mu_n, \mu)\to 0$ for $n\to\infty$, and let $f: X\to \R$ be a $1$-Lipschitz map. Then, we have 
    \begin{align*}
        \left|\int f \ d\mu_n - \int f \ d\mu\right| 
        &\leq \sup_{f \ 1\text{-Lip}} \int f \ d(\mu_n - \mu) \\
        &= \inf_{\pi\in Cpl(\mu_n, \mu)} \int d \ d\pi \quad \text{by Kantorovich-Rubenstein} \\
        &= \mathcal{W}_1(\mu_n, \mu) \leq \mathcal{W}_p(\mu_n, \mu) \to 0
    \end{align*}
    by Proposition \ref{prop:wasserstein_space_ordering}. The thesis follows by observing that on bounded Polish spaces $X$, the weak convergence in $\mathcal{P}(X)$ by continuous, bounded test functions is equivalent to the weak convergence by $1$-Lipschitz test functions.
\end{proof}

We will now remove the boundedness hypothesis on the previous theorem, and characterize completely the topology metrized by the Wasserstein distance.

\begin{definition}
    A sequence of measures $\{\mu_n\}_n\subseteq\mathcal{P}_p(X)$ converges weakly to a measure $\mu\in\mathcal{P}_p(X)$ if $\mu_n \to \mu$ weakly in $\mathcal{P}(X)$ and
    $$
    \int d(x, x_0)^p \ d\mu_n \underset{n\to\infty}{\longrightarrow} \int d(x, x_0)^p \ d\mu.
    $$
    for some $x_0\in X$.
\end{definition}

There exists many equivalent ways of introducing a notion of convergence in $\mathcal{P}_p(X)$.

\begin{proposition}
    \label{prop:equiv_statement_weak_conv_P_p(X)}
    Let $\{\mu_n\}_n\subseteq\mathcal{P}_p(X)$ and $\mu\in\mathcal{P}_p(X)$. Then, the following are equivalent:
    \begin{enumerate}[(i)]
        \item $\mu_n \to \mu$ weakly in $\mathcal{P}_p(X)$
        \item $\mu_n \to \mu$ weakly in $\mathcal{P}(X)$, and 
        $$
        \limsup_{n\to\infty} \int d(x, x)^p \ d\mu_n(x) \leq \int d(x, x_0)^p \ d\mu.
        $$
        \item $\mu_n\to \mu$ weakly in $\mathcal{P}(X)$ and
        $$
        \lim_{R\to\infty} \limsup_{n\to\infty} \int_{d(x, x_0)\geq R} d(x, x_0)^p \ d\mu(x) = 0.
        $$
        \item If $f: X\to \R$ continuous such that $|f(x)|\leq K(1+d(x, x_0)^p)$ for some $K\in\R$, then
        $$
        \int f \ d\mu_n \to \int f \ d\mu. 
        $$
    \end{enumerate}
\end{proposition}

\begin{lemma}
    \label{lemma:cauchy_in_P_p(X)}
    Let $\{\mu_n\}_n\subseteq \mathcal{P}_p(X)$ be a Cauchy sequence respect to $\mathcal{W}_p$. \\
    Then, $\{\mu_n\}_n$ is tight
\end{lemma}
\begin{proof}
    Lemma 6.14, \cite{villani2009optimal}.
\end{proof}

\begin{theorem}
    $\mathcal{W}_p$ metrizes the weak convergence in $\mathcal{P}_p(X)$, for $p\geq 1$, i.e. $$\mu_n\to\mu \text{ weakly in } \mathcal{P}_p(X) \iff \mathcal{W}_p(\mu_n, \mu) \to 0.$$
\end{theorem}
\begin{proof}
    Theorem 6.9, \cite{villani2009optimal}.
\end{proof}

\begin{theorem}
    Let $X$ be a Polish space, then $(\mathcal{P}_p(X), \mathcal{W}_p)$ Polish space for $p\geq 1$.
\end{theorem}
\begin{proof}
    Theorem 6.18, \cite{villani2009optimal}.
\end{proof}


\subsection{Displacement Interpolation}
In this section, we discuss the idea that leads to the dynamic formulation of optimal transport, namely the displacement interpolation. \\
Recall that, in $X=Y=\R^d$, there exists a unique constant speed geodesic between any couple of fixed points $x_0, x_1\in\R^d$ given by
$$
x_t := (1-t)x_0 + tx_1, \quad \text{for } t\in[0,1].
$$
Explicitly, $x_t$ is the unique optimizer of the problem
$$
\inf_{\gamma_0=x_0, \gamma_1=x_1}\int_0^1 |\dot \gamma_t|^2 \ dt.
$$
In the following, we will generalize this result to constant speed geodesics in $\mathcal{P}_p(\R^d)$.

\begin{proposition}
    Let $\mu_0, \mu_1\in\mathcal{P}_p(\R^d)$, and $X_0, X_1$ $\R^d$-valued variables with laws $\mu_0, \mu_1$. Then, the curve $\{\mu_t\}_{t\in[0,1]}\subseteq \mathcal{P}_p(\R^d)$ given by $$\mu_t \sim (1-t)X_0 + tX_1$$ is a constant speed geodesic in $\mathcal{P}_p(\R^d)$ if and only if $(X_0, X_1)$ is $\mathcal{W}_p$-optimal.
\end{proposition}
\begin{proof}
    Let $\mu_0, \mu_1\in\mathcal{P}_p(\R^d)$, and let $\pi\in Cpl(\mu_0, \mu_1)$ be a $\mathcal{W}_p$-optimal coupling, i.e. 
    $$
    \mathcal{W}_p(\mu_0, \mu_1)^p =  \int_{\R^d\times\R^d} |x-y|^p \ d\pi(x, y).
    $$
    Now, consider $X_0, X_1$ $\R^d$-valued random variables such that $(X_0, X_1)\sim \pi$ and, define for all $t\in[0,1]$
    $$
    X_t := (1-t)X_0 + tX_1,
    $$
    with law $\mu_t$. Then, we have by construction 
    \begin{align*}
        \mathcal{W}_p(\mu_s, \mu_t) 
        &= \left(\inf_{X\sim\mu_s, Y\sim\mu_t} \E[|X-Y|^p]\right)^{1/p} \\
        &\leq \E[|X_s - X_t|^p]^{1/p} \\
        &= \E[|s-t|^p \ |X_0 - X_1|^p]^{1/p} \\
        &= |t-s| \ \mathcal{W}_p(\mu_0, \mu_1),
    \end{align*}
    since $(X_0, X_1)$ optimal coupling for $\mathcal{W}_p(\mu_0, \mu_1)$. Then, the curve $\{\mu_t\}_{t\in[0,1]}$ is a constant speed geodesic between $\mu_0$ and $\mu_1$. \\
    Conversely, all the constant speed geodesics between $\mu_0, \mu_1$ are of this type.
\end{proof}

\begin{remark}
    In particular, the uniqueness of geodesic is not guaranteed in general, as the Wasserstein problem may might admit non-unique optimizers.
\end{remark}

\begin{example}
    The unique constant speed geodesics between $\delta_x, \delta_y \in \mathcal{P}_p(\R^d)$ is the curve $\{\delta_{x_t}\}_{t\in[0,1]}$ with
    $$
    x_t := (1-t)x + ty.
    $$
    Indeed, we have
    $$
    \mathcal{W}_p(\delta_{x_s}, \delta_{x_t}) = |x_s - x_t| = |t-s| \ |x - y|
    $$
    since the Wasserstein distance coincides with the underlying metric for Dirac measures. \\
    However, note that the curve obtained by convex interpolation $\mu_t := (1-t)\delta_{x} + t\delta_y$ is not a constant speed geodesics from $\delta_x$ to $\delta_y$. Hence, it holds
    $$
    \mathcal{W}_p(\mu_s, \mu_t) = |s-t|^{1/p} \ |x - y| > |t-s| |x-y| = |t-s| \ \mathcal{W}_p(\delta_x, \delta_y). 
    $$
    Intuitively, the constant speed geodesic $\delta_{x_t}$ moves from $\delta_x$ to $\delta_y$ by changing its support, while the convex interpolation $\mu_t$ goes from $\delta_x$ to $\delta_y$ by moving mass from $x$ to $y$.
\end{example}

As mentioned before, in the deterministic setting, the curve $x_t = (1-t)x_0 + tx_1$ is the unique optimizer of the problem
$$
\inf_{\gamma_0=x_0, \gamma_1=x_1}\int_0^1 |\dot \gamma_t|^2 \ dt.
$$
A similar result is guaranteed by the Brenier's theorem in the case of constant speed geodesics in $\mathcal{P}_2(\R^d)$.

\begin{proposition}
    The constant speed and absolutely continuous geodesics between $\mu, \nu\in \mathcal{P}_2(\R^d)$, with $\mu\ll\lambda$, are the laws of the optimizers of the problem 
    $$
    \mathcal{W}_2(\mu, \nu)^2 = \inf_{\substack{\{X_t\}_t\in AC(\R^d) \\ X_0 \sim\mu, \ X_1\sim\nu}} \int_0^1 \E[\dot X_t^2] \ dt,
    $$
    where $AC(\R^d)$ is the space of process with absolutely continuous paths $t\mapsto X_t\in\R^d$. 
\end{proposition}
\begin{proof}
    Let $\{\mu_t\}_t$ be a constant speed geodesic between $\mu, \nu$, i.e.
    $$
    \mu_t \sim X_t := (1-t)X_0 + tX_1 \quad \text{for } t\in[0,1],
    $$
    where $(X_0, X_1)$ is $\mathcal{W}_2(\mu, \nu)$-optimal, and $t\mapsto X_t$ is absolutely continuous. \\
    Now, note that the $2$-Wasserstein problem coincides with the Kantorovich problem with quadratic cost $c(x, y) = \frac{1}{2}|x-y|^2$, since
    $$
    \mathcal{W}_2(\mu, \nu)^2 = \inf_{\pi\in Cpl(\mu, \nu)} \int_{\R^d\times\R^d} |x-y|^2 \ d\pi(x, y) = \inf_{\pi\in Cpl(\mu, \nu)} \int_{\R^d\times\R^d} \frac{1}{2}|x-y|^2 \ d\pi(x, y).
    $$
    Then, there exists a unique optimal transport map from $\mu$ to $\nu$ of the form $\nabla\varphi: \R^d \to \R^d$, for some convex function $\varphi: \R^d\to\R$, by Brenier's theorem, i.e. it holds $$\nabla\varphi(\mu) = \nu.$$
    Equivalently, since $X_0\sim\mu$ and $X_1\sim\nu$, we also have $\nabla\varphi(X_0) = X_1$, which implies $X_t = (1-t)X_0 +t\nabla\varphi(X_0)$ with $$\dot X_t = \nabla\varphi(X_0) - X_0.$$
    Finally, we get
    $$
    \mathcal{W}_2(\mu, \nu)^2 = \E[|X_0 - X_1|^2] =\int_0^1 \E[\dot X_t^2] \ dt.
    $$
    Moreover, if $\{X_t\}_{t\in[0,1]}$ is an absolutely continuous curve with $X_0\sim\mu$ and $X_1\sim\nu$, then we have
    \begin{align*}
        \mathcal{W}_2(\mu, \nu)^2 
        &= \inf_{X\sim\mu,\ Y\sim\nu} \E[|X-Y|^2]
        \leq \E[|X_0 - X_1|^2] \\
        &\leq \E\left[\left|\int_0^1 \dot X_t \ dt\right|^2\right]
        \leq \E\left[\int_0^1 \dot X_t^2 \ dt\right] \quad \text{by Jensen}\\
        &= \int_0^1 \E[\dot X_t^2] \ dt \quad  \text{by Fubini}.
    \end{align*}
    The thesis follows by taking the infimum over all absolutely continuous curves with $X_0\sim\mu, \ X_1\sim\nu$.
\end{proof}

\subsection{Multi-marginal Optimal Transport}
In this section, we briefly discuss a multidimensional generalization of the optimal transport problem in the Kantorovich formulation, originally introduced in \cite{gangbo1998optimal}. We then present the Wasserstein barycenter problem as a particular instance of a multi-marginal optimal transport problem.

\begin{definition}
    A joint probability measure of $\mu_1, \dots, \mu_n\in\mathcal{P}(X)$ is a probability measure $\pi\in\mathcal{P}(X^n)$ with marginals $\mu_1, \dots, \mu_n$, i.e.
    $$
    p_i\pi = \mu_i \text{ for } i=1, \dots, n
    $$
    where $p_i: X^i \to X$ is the projection on the $i$-th component.
\end{definition}
We will denote the space of joint probabilities of $\mu_1, \dots, \mu_n$ by $\Pi(\mu_1, \dots, \mu_n)$. 

\begin{definition}
    A cost function for an optimal transport problem with marginals $\mu_1, \dots, \mu_n\in\mathcal{P}(X)$ is a measurable map $c: X^n\to \extRplus$ such that
    $$
    c(x_1, \dots, x_n)
    \geq
    \sum_{i=0}^n a_i(x_i) \text { for all } x_1, \dots, x_n\in X
    $$
    for some functions $a_1, \dots, a_n : X\to \extRminus$ with $a_i\in L^1(\mu_i)$ for $i=1, \dots, n$.
\end{definition}
Again, the lower bound assumption $c(x_1, \dots, x_n)\geq\sum_{i=0}^n a_i(x_i)$ guarantees that the expected cost $\E[c(X_1,\dots, X_n)]$ is well defined in $\extRplus$, with $X_i\sim \mu_i$.

\begin{definition}
    The $n$-marginal Kantorovich functional with cost $c: X^n \to \extRplus$ is the map $V_c : \mathcal{P}(X)^n \to \extRplus$ such that
    $$
    V^n_c(\mu_1, \dots, \mu_n) := \inf_{\pi\in \Pi(\mu_1, \dots, \mu_n)} \int_{X^n} c \ d\pi.
    $$
\end{definition}

\begin{definition}
    The $n$-marginal dual Kantorovich functional with cost $c: X^n\to \extRplus$ is the map $D^n_c(\mu_1, \dots, \mu_n) : \mathcal{P}(X)^n \to \extRplus$ such that 
    $$
    D^n_c(\mu_1, \dots, \mu_n) := \sup_{(\varphi_1, \dots, \varphi_n)\in D^n(c)} \sum_{i=1}^n \int_X \varphi_i \ d\mu_i
    $$
    with $D^n(c) := \{(\varphi_1, \dots, \varphi_n)\in L^1(\mu_1)\times\dots\times L^1(\mu_n) \mid c \geq \bigoplus_{i=1}^n \varphi_i\}$ research space.
\end{definition}

The classical existence and duality results for $2$-marginal optimal transport (Theorem \ref{thm:optimal_plan_existence}, Theorem \ref{thm:strong_duality}) can be adapted to the $n$-dimensional setting.

\begin{theorem}
\label{thm:multi_dim_fund_thm_ot}
    Let $c: X^n \to \R$ be a LSC cost function for a $n$-dimensional optimal transport problem with marginals $\mu_1, \dots, \mu_n\in\mathcal{P}(X)$. Then, the following holds:
    \begin{enumerate}[(i)]
        \item If $c\geq 0$, then the map $$\pi\mapsto \int_{X^n} c \ d\pi$$ is lower-semicontinuous.
        \item There exists a minimizer for the $n$-marginal Kantorovich problem, i.e. $$\exists \ \pi^*\in\argmin_{\pi\in\Pi(\mu_1, \dots, \mu_n)} \int_{X^n} c \ d\pi.$$
        \item Duality holds, i.e. $D^n_c(\mu_1, \dots, \mu_n) = V^n_c(\mu_1, \dots, \mu_n)$.
    \end{enumerate}
\end{theorem}

\begin{remark}
    The geometry of optimal joint measures becomes substantially richer for $n\geq 3$, and understanding their structure remains an active area of research. Indeed, uniqueness of optimizers and concentration on the graph of a transport map are no longer typical properties of multi-dimensional optimal transport problems.
\end{remark}

A particular instance of a multi-marginal Kantorovich problem is given by the Wasserstein barycenter problem, firstly introduced in 2011 by \cite{Agueh2011bayecenter}.

\begin{definition}
    The $p$-Wasserstein barycenter functional is the map $V^{W}: \mathcal{P}_p(\R^d)^N \to \extRplus$ such that
    $$
    V^W(\mu_1, \dots, \mu_n) := \inf_{\mu\in\mathcal{P}_p(\R^d)} \frac{1}{N}\sum_{i=1}^N \mathcal{W}_p(\mu_i, \mu)^p.
    $$
    In particular, a Wasserstein barycenter of $\mu_1, \dots, \mu_n \in\mathcal{P}_p(\R^d)$ is a minimizer of the optimization problem.
\end{definition}

\begin{remark}
    Intuitively, a Wasserstein barycenter of $\mu_1, \dots, \mu_N\in\mathcal{P}_p(X)$ is a notion of mean in the Wasserstein space $(\mathcal{P}_p(\R^d), \mathcal{W}_p)$. Moreover, in general the Wasserstein barycenter is not unique, reflecting the geometric complexity of $\mathcal{P}_p(X)$.
\end{remark}

The next result shows that we can regard the $2$-Wasserstein barycenter problem as a multi-marginal Kantorovich problem for a specific quadratic cost function.

\begin{theorem}
    Let $\mu_1, \dots, \mu_n\in\mathcal{P}_2(\R^d)$ be probability measures with finite second moment. Then, the $2$-Wasserstein barycenter problem is a multi-marginal Kantorovich problem with cost $c : {(\R^d)}^N \to \R$ such that $c(x_1, \dots, x_N) := \frac{1}{N} \sum_{i=1}^N \| x_i - T(x_1, \dots, x_n)\| ^2$, i.e.
    $$
    V^W(\mu_1, \dots, \mu_n) := \inf_{\mu\in\mathcal{P}_2(\R^d)} \frac{1}{N}\sum_{i=1}^N \mathcal{W}_2(\mu_i, \mu)^2
    = V^N_c(\mu_1, \dots, \mu_N).
    $$
    where $T(x_1, \dots, \mu_N) = \bar x_N = \frac{1}{N} \sum_{i=1}^N x_i \in \R^d$ average point. \\
    Moreover, if $\pi^*$ is an optimizer of $V^N_c(\mu_1, \dots, \mu_N)$, then $T\pi^*$ is a $2$-Wasserstein barycenter of $\mu_1, \dots, \mu_N$.
\end{theorem}
\begin{proof}
    We have
    \begin{align*}
        \inf_{\mu\in\mathcal{P}_2(\R^d)} \frac{1}{N}\sum_{i=1}^N \mathcal{W}_2(\mu_i, \mu)^2
        &=  \inf_{\mu\in\mathcal{P}_2(\R^d)} \frac{1}{N}\sum_{i=1}^N \inf_{X_i\sim\mu_i, X\sim\mu} \E[\|X_i - X\|^2] = \\
        &= \inf_{X_1\sim\mu_1, \dots, X_n\sim\mu_N} \ \inf_{\mu\in\mathcal{P}_2(\R^d), X\ \sim \mu} \  \frac{1}{N}\sum_{i=1}^N \E[\|X_i - X\|^2] = \\
        &= \inf_{X_1\sim\mu_1, \dots, X_n\sim\mu_N} \ \frac{1}{N}\sum_{i=1}^N \E[\|X_i - \bar X_N\|^2]
    \end{align*}
    since the sample mean of the data points minimizes the squared loss. Then, we obtain
    \begin{align*}
        \inf_{\mu\in\mathcal{P}_2(\R^d)} \frac{1}{N}\sum_{i=1}^N \mathcal{W}_2(\mu_i, \mu)^2
        &= \inf_{X_1\sim\mu_1, \dots, X_n\sim\mu_N} \E\left[\frac{1}{N}\sum_{i=1}^N \|X_i - \bar X_N\|^2\right] \\
        &= \inf_{\pi\in\Pi(\mu_1, \dots, \mu_n)} \int_{(\R^d)^N} \frac{1}{N}\sum_{i=1}^N \|x_i - \bar x_N\|^2 \ d\pi(x_1, \dots, x_N) \\
        &= \inf_{\pi\in\Pi(\mu_1, \dots, \mu_n)} \int_{(\R^d)^N} c(x_1, \dots, x_N) \ d\pi(x_1, \dots, x_N) \\
        &= V^n_c(\mu_1, \dots, \mu_N).
    \end{align*}
    Thus, $W(\mu_1, \dots, \mu_N)$ is a multi-marginal Kantorovich problem with quadratic cost $c$. Moreover, by Theorem \ref{thm:multi_dim_fund_thm_ot} the infimum is attained by some optimizer $\pi^*$ and the induced measure $T\pi^*$ on $\R^d$ by pushforward under the map $T: (\R^d)^N\to \R^d$, is a $2$-Wasserstein barycenter of $\mu_1,\dots, \mu_N$.
\end{proof}
\end{document}